# Hello, world!
#
# This is an example function named 'hello'
# which prints 'Hello, world!'.
#
# You can learn more about package authoring with RStudio at:
#
#   https://r-pkgs.org
#
# Some useful keyboard shortcuts for package authoring:
#
#   Install Package:           'Ctrl + Shift + B'
#   Check Package:             'Ctrl + Shift + E'
#   Test Package:              'Ctrl + Shift + T'
# install.packages("xtable")
# Install stargazer if it's not already installed
# install.packages("xtable")
# Install stargazer if it's not already installed
#  (!require("stargazer")) {
# install.packages("stargazer")
#

#  (!requireNamespace("gsynth", quietly = TRUE)) {
# install.packages("gsynth")
#


# install.packages("MatchIt")
# install.packages("boot")
# install.packages("Matching")
# install.packages("cobalt")
# install.packages("plm")
library(plm)
library(boot)
library(Matching)
library(dplyr)
library(tidyr)
library(zoo)
library(readxl)
library(MatchIt)
library(stargazer)  # Load stargazer for summary statistics
library(ggplot2)
library(xtable)
library(cobalt)
library(kableExtra)
library(texreg)
library(gsynth)
library(WeightIt)
library(ragg)




# Set working directories for saving tables and figures
tables_dir <- "C:/Users/ammonsj/Dropbox/Apps/Overleaf/Gender and Economic Freedom/Tables"  # [INSERT]
figures_dir <- "C:/Users/ammonsj/Dropbox/Apps/Overleaf/Gender and Economic Freedom/Figures"  # [INSERT]

# Ensure directories exist
dir.create(tables_dir, showWarnings = FALSE, recursive = TRUE)
dir.create(figures_dir, showWarnings = FALSE, recursive = TRUE)

# Function to save tables
save_table_results <- function(tab, caption, file_name) {
  stargazer(tab,
            type = "latex",
            title = caption,
            out = file_name)
}


pwt1001 <- read_excel("C:/Users/ammonsj/Downloads/pwt1001.xlsx",
                      sheet = "Data")
View(pwt1001)

P_Data_Extract_From_World_Development_Indicators_4_ <- read_excel("C:/Users/ammonsj/Downloads/P_Data_Extract_From_World_Development_Indicators (4).xlsx")
View(P_Data_Extract_From_World_Development_Indicators_4_)

efotw_2023_master_index_data_for_researchers_iso_1_ <- read_excel("C:/Users/ammonsj/Downloads/efotw-2023-master-index-data-for-researchers-iso (1).xlsx")
View(efotw_2023_master_index_data_for_researchers_iso_1_)


# Create the dataset from efotw_2023_master_index_data_for_researchers_iso_1_
efotw <- efotw_2023_master_index_data_for_researchers_iso_1_

# Adjust the numeric_columns list based on the actual column names
numeric_columns <- c("Economic Freedom Summary Index", "Rank", "Quartile",
                     "1A Government consumption", "1B  Transfers and subsidies",
                     "1C  Government investment", "1Di Top marginal income tax rate",
                     "1Dii Top marginal income and payroll tax rate", "1D  Top marginal tax rate",
                     "IE State ownership", "1  Size of Government", "Area 1 Rank",
                     "2A  Judicial independence", "2B  Impartial courts", "2C  Property rights",
                     "2D  Military interference", "2E Legal integrity", "2F Contracts",
                     "2G Real property", "2H Police and crime", "Gender Disparity Index",
                     "2  Legal System & Property Rights -- With Gender Adjustment", "Area 2 Rank",
                     "2 Legal System & Property Rights - No Gender Adjustment", "3A  Money growth",
                     "3B  Standard deviation of inflation", "3C  Inflation",
                     "3D  Foreign currency bank accounts", "3  Sound Money", "Area 3 Rank",
                     "4Ai  Trade tax revenue", "4Aii  Mean tariff rate", "4Aiii  Standard deviation of tariff rates",
                     "4A  Tariffs", "4Bi  Non-tariff trade barriers", "4Bii  Costs of importing and exporting",
                     "4B  Regulatory trade barriers", "4C  Black market exchange rates",
                     "4Di  Financial openness", "4Dii  Capital controls", "4Diii Freedom of foreigners to visit",
                     "4Div Protection of Foreign Assets", "4D  Controls of the movement of capital and people",
                     "4  Freedom to trade internationally", "Area 4 Rank",
                     "5Ai  Ownership of banks", "5Aii Private sector credit",
                     "5Aiii  Interest rate controls/negative real interest rates)", "5A  Credit market regulation",
                     "5Bi  Labor regulations and minimum wage", "5Bii  Hiring and firing regulations",
                     "5Biii  Flexible wage determination", "5Biv  Hours Regulations", "5Bv Cost of worker dismissal",
                     "5Bvi  Conscription", "5Bvii Foreign Labor", "5B  Labor market regulations",
                     "5Ci  Regulatory Burden", "5Cii  Bureacracy costs", "5Ciii  Impartial Public Administration",
                     "5Civ Tax compliance", "5C  Business regulations", "5Di  Market openness",
                     "5Dii Business Permits", "5D Freedom to enter markets and compete", "5  Regulation", "Area 5 Rank")

`V-Dem-CY-Full+Others-v15` <- readRDS("C:/Users/ammonsj/OneDrive - Wabash College/Desktop/Databases/V-Dem-CY-Full+Others-v15.rds")

# For pwt1001
colnames(pwt1001)

# For P_Data_Extract_From_World_Development_Indicators_4_
colnames(P_Data_Extract_From_World_Development_Indicators_4_)

# For efotw_2023_master_index_data_for_researchers_iso_1_
colnames(efotw_2023_master_index_data_for_researchers_iso_1_)


# Load the datasets
pwt <- read_excel("C:/Users/ammonsj/Downloads/pwt1001.xlsx", sheet = "Data")
# Load the World Development Indicators dataset while skipping non-data rows
wdi <- read_excel("C:/Users/ammonsj/Downloads/P_Data_Extract_From_World_Development_Indicators (4).xlsx")
vdem <- readRDS("C:/Users/ammonsj/OneDrive - Wabash College/Desktop/Databases/V-Dem-CY-Full+Others-v15.rds")

# Renaming columns in wdi and efotw
wdi <- rename(wdi, countrycode = `Country Code`, year = Time)
efotw <- rename(efotw, countrycode = `ISO Code 3`, year = `Year`)

# Assuming 'country_text_id' in 'vdem' is the correct country code column
# If 'country_text_id' contains ISO country codes; otherwise, check the appropriate column
vdem <- rename(vdem, countrycode = country_text_id)

vdem <- vdem %>%
  select(countrycode, year, v2x_polyarchy, e_peaveduc, e_peedgini, v2clprptyw)

# Filter out data before 1970 and after 2021 for each dataset
pwt <- pwt %>% filter(year >= 1970 & year <= 2021)
wdi <- wdi %>% filter(year >= 1970 & year <= 2021)
efotw <- efotw %>% filter(year >= 1970 & year <= 2021)
vdem <- vdem %>% filter(year >= 1970 & year <= 2021)

# Now attempt to merge again
merged_data <- pwt %>%
  left_join(wdi, by = c("countrycode", "year")) %>%
  left_join(vdem, by = c("countrycode", "year")) %>%
  left_join(efotw, by = c("countrycode", "year"))


# Review the structure of the merged dataset to ensure all variables are correctly aligned
str(merged_data)

# View the filtered dataset
View(merged_data)

merged_data <- merged_data %>%
  rename(
    school_enrollment_primary = `School enrollment, primary (% net) [SE.PRM.NENR]`,
    school_enrollment_secondary = `School enrollment, secondary (% net) [SE.SEC.NENR]`,
    maternal_mortality = `Maternal mortality ratio (modeled estimate, per 100,000 live births) [SH.STA.MMRT]`,
    sanitation_urban = `People using at least basic sanitation services, urban (% of urban population) [SH.STA.BASS.UR.ZS]`,
    life_expectancy = `Life expectancy at birth, total (years) [SP.DYN.LE00.IN]`,
    unemployment_rate = `Unemployment, total (% of total labor force) (modeled ILO estimate) [SL.UEM.TOTL.ZS]`,
    adj_net_income_per_capita = `Adjusted net national income per capita (current US$) [NY.ADJ.NNTY.PC.CD]`,
    govt_consumption_expenditure = `General government final consumption expenditure (% of GDP) [NE.CON.GOVT.ZS]`,
    exports_gdp = `Exports of goods and services (% of GDP) [NE.EXP.GNFS.ZS]`,
    inflation = `Inflation, consumer prices (annual %) [FP.CPI.TOTL.ZG]`
  )


# Replace ".." with NA in character columns only
merged_data <- merged_data %>%
  mutate(across(where(is.character), ~na_if(.x, "..")))
# Replace spaces with underscores in all column names
merged_data <- merged_data %>%
  rename_with(~ gsub(" ", "_", .x))

# Filter out all rows where year is before 1970
merged_data <- merged_data %>% filter(year >= 1970)

# Convert relevant columns in merged_data to numeric
merged_data$school_enrollment_primary <- suppressWarnings(as.numeric(merged_data$school_enrollment_primary))
merged_data$school_enrollment_secondary <- suppressWarnings(as.numeric(merged_data$school_enrollment_secondary))
merged_data$maternal_mortality <- suppressWarnings(as.numeric(merged_data$maternal_mortality))
merged_data$sanitation_urban <- suppressWarnings(as.numeric(merged_data$sanitation_urban))
merged_data$life_expectancy <- suppressWarnings(as.numeric(merged_data$life_expectancy))

# Optionally, check if conversion was successful
str(merged_data)


columns_to_keep <- c("countrycode", "year", "Gender_Disparity_Index", "inflation",
                     "rgdpe", "pop", "csh_g", "csh_x", "v2x_polyarchy",
                     "hc", "school_enrollment_primary", "school_enrollment_secondary",
                     "maternal_mortality", "sanitation_urban", "life_expectancy",
                     "unemployment_rate", "adj_net_income_per_capita",
                     "govt_consumption_expenditure", "exports_gdp",
                     "Economic_Freedom_Summary_Index", "e_peaveduc",
                     "e_peedgini", "efw", "v2clprptyw")


# Retain only the specified columns
merged_data <- merged_data %>%
  select(all_of(columns_to_keep))

# Complete missing years and fill NA gaps using interpolation
merged_data <- merged_data %>%
  group_by(countrycode) %>%
  complete(year = full_seq(year, 1)) %>%  # Fill in missing years
  mutate(across(where(is.numeric), ~ na.approx(., na.rm = FALSE)))  # Interpolate numeric columns

merged_data <- merged_data %>%
  mutate(rgdpe_per_capita = rgdpe / pop)


# Step 1: Sort data by countrycode and year
merged_data <- merged_data %>%
  arrange(countrycode, year)

# Step 2: Calculate 5-year cumulative change in GDI for each country
merged_data <- merged_data %>%
  group_by(countrycode) %>%
  mutate(GDI_change_5yr = Gender_Disparity_Index - lag(Gender_Disparity_Index, 5))

merged_data <- merged_data %>%
  arrange(countrycode, year) %>%  # Ensure data is sorted correctly by country and year
  group_by(countrycode) %>%
  mutate(
    # Calculate the rolling 5-year change in Gender Disparity Index
    GDI_change_5yr = Gender_Disparity_Index - lag(Gender_Disparity_Index, 5),
    # Flag the treatment as 1 if the 5-year GDI change is greater than or equal to 0.10
    treatment_5yr = ifelse(GDI_change_5yr >= 0.15, 1, 0)  ) %>%
  ungroup()

merged_data <- merged_data %>%
  group_by(countrycode) %>%
  mutate(GDI_change_5yr.125 = Gender_Disparity_Index - lag(Gender_Disparity_Index, 5))

merged_data <- merged_data %>%
  arrange(countrycode, year) %>%  # Ensure data is sorted correctly by country and year
  group_by(countrycode) %>%
  mutate(
    # Calculate the rolling 5-year change in Gender Disparity Index
    GDI_change_5yr.125 = Gender_Disparity_Index - lag(Gender_Disparity_Index, 5),
    # Flag the treatment as 1 if the 5-year GDI change is greater than or equal to 0.10
    treatment_5yr.125 = ifelse(GDI_change_5yr >= 0.125, 1, 0)  ) %>%
  ungroup()


merged_data <- merged_data %>%
  group_by(countrycode) %>%
  mutate(GDI_change_5yr.1 = Gender_Disparity_Index - lag(Gender_Disparity_Index, 5))

merged_data <- merged_data %>%
  arrange(countrycode, year) %>%  # Ensure data is sorted correctly by country and year
  group_by(countrycode) %>%
  mutate(
    # Calculate the rolling 5-year change in Gender Disparity Index
    GDI_change_5yr.1 = Gender_Disparity_Index - lag(Gender_Disparity_Index, 5),
    # Flag the treatment as 1 if the 5-year GDI change is greater than or equal to 0.10
    treatment_5yr.1 = ifelse(GDI_change_5yr >= 0.1, 1, 0)  ) %>%
  ungroup()

#==============================================================================
# PART 2: MASTER ANALYSIS FUNCTION DEFINITION (REVISED FOR FORMATTING)
#==============================================================================
# This single, robust function performs the entire analysis for any given
# threshold. It creates and saves all plots and tables.

run_gdi_jump_analysis <- function(data, jump_threshold, figures_dir, tables_dir) {

  file_suffix <- gsub("\\.", "", as.character(jump_threshold))
  cat(paste("\n\n===== Running Analysis for GDI Jump Threshold: >=", jump_threshold, "=====\n"))

  # --- 1. Event Study Setup (No Changes) ---
  analysis_data <- data %>%
    mutate(treatment_var = ifelse(!is.na(GDI_change_5yr) & GDI_change_5yr >= jump_threshold, 1, 0))

  treatment_events <- analysis_data %>%
    filter(treatment_var == 1) %>%
    group_by(countrycode) %>%
    summarise(treatment_year = min(year), .groups = 'drop')

  if (nrow(treatment_events) == 0) {
    cat(paste(">>> WARNING: No treatment events found for threshold", jump_threshold, ". Skipping analysis.\n"))
    return(invisible(NULL))
  }

  analysis_data <- analysis_data %>%
    left_join(treatment_events, by = "countrycode") %>%
    mutate(time_to_treatment = ifelse(!is.na(treatment_year), year - treatment_year, NA_integer_)) %>%
    group_by(countrycode) %>%
    mutate(prop_rights_change = v2clprptyw - lag(v2clprptyw, 1)) %>%
    ungroup()

  # --- 2. Time-Series Line Plot (Formatting Revised) ---
  time_series_data <- analysis_data %>%
    filter(!is.na(time_to_treatment) & time_to_treatment >= -20 & time_to_treatment <= 20) %>%
    filter(!is.na(prop_rights_change)) %>%
    group_by(time_to_treatment) %>%
    summarise(
      mean_change = mean(prop_rights_change, na.rm = TRUE),
      se = sd(prop_rights_change, na.rm = TRUE) / sqrt(n()),
      lower_ci = mean_change - 1.96 * se,
      upper_ci = mean_change + 1.96 * se,
      .groups = 'drop'
    )

  if (nrow(time_series_data) > 0) {
    max_y_val <- max(time_series_data$upper_ci, na.rm = TRUE)
    
    # === PLOT FORMATTING CHANGES START HERE ===
    time_series_plot <- ggplot(time_series_data, aes(x = time_to_treatment, y = mean_change)) +
      geom_ribbon(aes(ymin = lower_ci, ymax = upper_ci), alpha = 0.3, fill = "gray60") +
      geom_line(color = "black", linewidth = 1.5) +
      geom_point(color = "black", size = 3) +
      geom_hline(yintercept = 0, linetype = "solid", color = "gray30") +
      geom_vline(xintercept = 0, linetype = "dashed", color = "black", linewidth = 1.2) +
      geom_vline(xintercept = -5, linetype = "dotted", color = "black", linewidth = 1.2) +
      scale_x_continuous(breaks = seq(-20, 20, 5)) +
      theme_bw() + # Changed to a black & white theme
      theme(     # Increased all font sizes and weights
        panel.grid.minor = element_blank(),
        plot.title = element_text(size = 20, face = "bold", hjust = 0.5),
        plot.subtitle = element_text(size = 16, hjust = 0.5),
        axis.title = element_text(size = 16, face = "bold"),
        axis.text = element_text(size = 14, color = "black")
      ) +
      labs(
        title = paste("Evolution of Property Rights Around GDI Jumps (Threshold ≥", jump_threshold, ")"),
        subtitle = "Mean year-over-year change with 95% confidence intervals",
        x = paste("Years Relative to GDI Jump"),
        y = "Mean Annual Change in Property Rights for Women"
      ) +
      # Annotations are now black and larger
      annotate("text", x = 0, y = max_y_val * 1.1, label = "GDI Jump", color = "black", fontface = "bold", size = 5.5) +
      annotate("text", x = -5, y = max_y_val * 1.1, label = "Jump Period\nBegins", color = "black", fontface = "bold", size = 5.5, hjust = 0.5)
    # === PLOT FORMATTING CHANGES END HERE ===

    ggsave(file.path(figures_dir, paste0("property_rights_evolution_line_plot_jump_", file_suffix, ".png")), time_series_plot, width = 12, height = 7, dpi = 1200)
    print(time_series_plot)
  }

  # --- 3. Chronologically Ordered Summary Table (No Changes) ---
  period_summary <- analysis_data %>%
    mutate(period = factor(
      case_when(
        time_to_treatment >= -20 & time_to_treatment < -5  ~ "Pre-Jump (-20 to -5)",
        time_to_treatment >= -5 & time_to_treatment <= 0  ~ "During Jump (-5 to 0)",
        time_to_treatment > 0 & time_to_treatment <= 20   ~ "Post-Jump (0 to +20)"
      ),
      levels = c("Pre-Jump (-20 to -5)", "During Jump (-5 to 0)", "Post-Jump (0 to +20)")
    )) %>%
    filter(!is.na(period) & !is.na(prop_rights_change)) %>%
    group_by(period) %>%
    summarise(
      `Mean Annual Change` = round(mean(prop_rights_change, na.rm = TRUE), 4),
      `Std. Error` = round(sd(prop_rights_change, na.rm = TRUE) / sqrt(n()), 4),
      `95% CI Lower` = round(`Mean Annual Change` - 1.96 * `Std. Error`, 4),
      `95% CI Upper` = round(`Mean Annual Change` + 1.96 * `Std. Error`, 4),
      `N (Country-Years)` = n(),
      .groups = 'drop'
    ) %>%
    arrange(period)

  if (nrow(period_summary) > 0) {
    stargazer(as.data.frame(period_summary), type = "latex", summary = FALSE,
              title = paste("Property Rights Changes Around GDI Jumps (Threshold:", jump_threshold, ")"),
              label = paste0("tab:prop_rights_summary_jump_", file_suffix),
              out = file.path(tables_dir, paste0("property_rights_summary_table_jump_", file_suffix, ".tex")))
    cat(paste("\n--- Summary Table for Threshold:", jump_threshold, "---\n"))
    print(period_summary)
  }

  # --- 4. Chronological Bar Chart (Formatting Revised) ---
  if (nrow(period_summary) > 0) {
    
    # === PLOT FORMATTING CHANGES START HERE ===
    summary_bar_chart <- ggplot(period_summary, aes(x = period, y = `Mean Annual Change`)) +
      geom_col(fill = "gray50", color = "black") + # Gray bars with black outline
      geom_errorbar(aes(ymin = `95% CI Lower`, ymax = `95% CI Upper`), width = 0.25, color = "black", linewidth = 1) + # Thicker error bars
      geom_hline(yintercept = 0, color = "gray30") + # Reference line at zero
      labs(
        title = paste("Mean Change in Women's Property Rights (Jump ≥", jump_threshold, ")"),
        subtitle = "Periods are shown in chronological order",
        x = "Period Relative to GDI Jump",
        y = "Mean Annual Change in Property Rights"
      ) +
      theme_bw() + # Changed to a black & white theme
      theme(     # Increased all font sizes and weights
        panel.grid.major.x = element_blank(),
        panel.grid.minor = element_blank(),
        plot.title = element_text(size = 20, face = "bold", hjust = 0.5),
        plot.subtitle = element_text(size = 16, hjust = 0.5),
        axis.title = element_text(size = 16, face = "bold"),
        axis.text = element_text(size = 14, color = "black")
      )
    # === PLOT FORMATTING CHANGES END HERE ===

    ggsave(file.path(figures_dir, paste0("property_rights_summary_barchart_jump_", file_suffix, ".png")), summary_bar_chart, width = 8, height = 6, dpi = 1200)
    print(summary_bar_chart)
  }
  cat(paste("\n===== Finished Analysis for Threshold:", jump_threshold, "=====\n"))
}


#==============================================================================
# PART 3: EXECUTE ANALYSIS FOR ALL THRESHOLDS
#==============================================================================
# This section calls the master function for each desired threshold.

# Run for the primary 0.15 jump threshold
run_gdi_jump_analysis(data = merged_data,
                      jump_threshold = 0.15,
                      figures_dir = figures_dir,
                      tables_dir = tables_dir)

# Run for the 0.125 robustness check
run_gdi_jump_analysis(data = merged_data,
                      jump_threshold = 0.125,
                      figures_dir = figures_dir,
                      tables_dir = tables_dir)

# Run for the 0.1 robustness check
run_gdi_jump_analysis(data = merged_data,
                      jump_threshold = 0.1,
                      figures_dir = figures_dir,
                      tables_dir = tables_dir)




# Load necessary packages
library(stargazer)
library(plm)
library(dplyr)

# Set up paths for saving output tables
tables_dir <- "C:/Users/ammonsj/Dropbox/Apps/Overleaf/Gender and Economic Freedom/Tables"

# Prepare the dataset
data <- merged_data

# List of outcome variables
outcome_vars <- c("rgdpe_per_capita", "maternal_mortality", "sanitation_urban", "life_expectancy")

# List of treatment variables corresponding to different thresholds
treatment_vars <- c("treatment_5yr.1", "treatment_5yr.125", "treatment_5yr")
treatment_labels <- c("0.10", "0.125", "0.15")

### Enhanced function to run models and output formatted LaTeX table ###
run_models_and_save_formatted_table <- function(outcome, treatment_var, threshold_label, data) {
  
  # Create safe file name
  file_suffix <- gsub("\\.", "", threshold_label)
  outcome_clean <- gsub("_", "", outcome)
  
  # Model 1: OLS (no fixed effects)
  model1 <- lm(as.formula(paste(outcome, "~", treatment_var)), data = data)
  
  # Model 2: Country Fixed Effects only
  model2 <- plm(as.formula(paste(outcome, "~", treatment_var)), 
                data = data, 
                index = c("countrycode", "year"), 
                model = "within", 
                effect = "individual")
  
  # Model 3: Year Fixed Effects only
  model3 <- plm(as.formula(paste(outcome, "~", treatment_var)), 
                data = data, 
                index = c("countrycode", "year"), 
                model = "within", 
                effect = "time")
  
  # Model 4: Two-way Fixed Effects (Country and Year)
  model4 <- plm(as.formula(paste(outcome, "~", treatment_var)), 
                data = data, 
                index = c("countrycode", "year"), 
                model = "within", 
                effect = "twoways")
  
  # Extract R-squared values
  r2_ols <- summary(model1)$r.squared
  r2_country <- summary(model2)$r.squared["rsq"]
  r2_year <- summary(model3)$r.squared["rsq"]
  r2_twoway <- summary(model4)$r.squared["rsq"]
  
  # Extract within R-squared values (for FE models)
  within_r2_country <- summary(model2)$r.squared["rsq"]
  within_r2_year <- summary(model3)$r.squared["rsq"]
  within_r2_twoway <- summary(model4)$r.squared["rsq"]
  
  # Create custom table title
  outcome_label <- switch(outcome,
    "rgdpe_per_capita" = "GDP per Capita",
    "maternal_mortality" = "Maternal Mortality",
    "sanitation_urban" = "Urban Sanitation",
    "life_expectancy" = "Life Expectancy",
    outcome  # Default to variable name if not in list
  )
  
  table_title <- paste0("Impact of GDI Jump ($\geq$", threshold_label, ") on ", outcome_label)
  
  # Create the LaTeX table with enhanced formatting
  stargazer(model1, model2, model3, model4,
            type = "latex",
            title = table_title,
            label = paste0("tab:", outcome_clean, "_", file_suffix),
            dep.var.labels.include = FALSE,
            column.labels = c("OLS", "Country FE", "Year FE", "Country + Year FE"),
            covariate.labels = c(paste0("GDI Jump ($\geq$", threshold_label, ")")),
            keep.stat = c("n", "rsq"),
            star.cutoffs = c(0.1, 0.05, 0.01),
            star.char = c("*", "**", "***"),
            notes = c("Standard errors in parentheses.", 
                     "* p$<$0.1, ** p$<$0.05, *** p$<$0.01"),
            notes.align = "l",
            add.lines = list(
              c("Country Fixed Effects", "No", "Yes", "No", "Yes"),
              c("Year Fixed Effects", "No", "No", "Yes", "Yes"),
              c("Observations", 
                format(length(model1$residuals), big.mark = ","),
                format(length(model2$residuals), big.mark = ","),
                format(length(model3$residuals), big.mark = ","),
                format(length(model4$residuals), big.mark = ",")),
              c("R²", 
                sprintf("%.3f", r2_ols),
                sprintf("%.3f", r2_country),
                sprintf("%.3f", r2_year),
                sprintf("%.3f", r2_twoway)),
              c("Within R²", 
                "—",
                sprintf("%.3f", within_r2_country),
                sprintf("%.3f", within_r2_year),
                sprintf("%.3f", within_r2_twoway))
            ),
            out = file.path(tables_dir, paste0("regression_", outcome_clean, "_jump", file_suffix, ".tex")))
  
  # Print confirmation
  cat(paste0("Created table for ", outcome_label, " with threshold $\geq$", threshold_label, "\n"))
}

### Master function to run all combinations ###
run_all_regression_tables <- function() {
  cat("\n========== Creating Regression Tables ==========\n\n")
  
  # Loop through each outcome variable
  for (i in seq_along(outcome_vars)) {
    outcome <- outcome_vars[i]
    cat(paste0("\n--- Processing outcome: ", outcome, " ---\n"))
    
    # Loop through each treatment threshold
    for (j in seq_along(treatment_vars)) {
      treatment <- treatment_vars[j]
      threshold <- treatment_labels[j]
      
      # Check if treatment variable exists in the data
      if (!(treatment %in% names(data))) {
        cat(paste0("Warning: Treatment variable '", treatment, "' not found in data. Skipping.\n"))
        next
      }
      
      # Check if outcome variable exists in the data
      if (!(outcome %in% names(data))) {
        cat(paste0("Warning: Outcome variable '", outcome, "' not found in data. Skipping.\n"))
        next
      }
      
      # Run the regression and create table
      tryCatch({
        run_models_and_save_formatted_table(outcome, treatment, threshold, data)
      }, error = function(e) {
        cat(paste0("Error creating table for ", outcome, " with threshold ", threshold, ": ", e$message, "\n"))
      })
    }
  }
  
  cat("\n========== All Regression Tables Created ==========\n")
}

# Execute the function to create all tables
run_all_regression_tables()

### Optional: Create a combined summary table for key outcomes ###
create_summary_comparison_table <- function(outcomes_subset = c("rgdpe_per_capita", "life_expectancy", "maternal_mortality")) {
  
  cat("\n========== Creating Summary Comparison Table ==========\n")
  
  # Store results for the summary table
  summary_results <- list()
  
  for (outcome in outcomes_subset) {
    for (j in seq_along(treatment_vars)) {
      treatment <- treatment_vars[j]
      threshold <- treatment_labels[j]
      
      # Skip if variables don't exist
      if (!(treatment %in% names(data)) || !(outcome %in% names(data))) {
        next
      }
      
      # Run two-way fixed effects model
      model <- plm(as.formula(paste(outcome, "~", treatment)), 
                   data = data, 
                   index = c("countrycode", "year"), 
                   model = "within", 
                   effect = "twoways")
      
      # Extract coefficient and standard error
      coef_val <- coef(summary(model))[1, "Estimate"]
      se_val <- coef(summary(model))[1, "Std. Error"]
      p_val <- coef(summary(model))[1, "Pr(>|t|)"]
      
      # Determine significance stars
      stars <- ""
      if (!is.na(p_val)) {
        if (p_val < 0.01) stars <- "***"
        else if (p_val < 0.05) stars <- "**"
        else if (p_val < 0.1) stars <- "*"
      }
      
      # Format the result
      result_str <- paste0(sprintf("%.3f", coef_val), stars, "\n(", sprintf("%.3f", se_val), ")")
      
      # Store result
      key <- paste0(outcome, "_", threshold)
      summary_results[[key]] <- result_str
    }
  }
  
  # Create a formatted comparison table (you can customize this further)
  cat("\n--- Two-Way Fixed Effects Results Summary ---\n")
  cat("Outcome | Jump $\geq$0.10 | Jump $\geq$0.125 | Jump $\geq$0.15\n")
  cat("--------|------------|-------------|------------\n")
  
  for (outcome in outcomes_subset) {
    outcome_label <- switch(outcome,
      "rgdpe_per_capita" = "GDP per Capita",
      "life_expectancy" = "Life Expectancy",
      "maternal_mortality" = "Maternal Mortality",
      outcome
    )
    
    row_text <- paste0(outcome_label, " | ")
    for (threshold in treatment_labels) {
      key <- paste0(outcome, "_", threshold)
      if (key %in% names(summary_results)) {
        row_text <- paste0(row_text, summary_results[[key]], " | ")
      } else {
        row_text <- paste0(row_text, "N/A | ")
      }
    }
    cat(row_text, "\n")
  }
  
  cat("\n========== Summary Table Complete ==========\n")
}

# Optional: Create the summary comparison table for key outcomes
create_summary_comparison_table()






# Step 5: Create lagged variables for propensity score matching
merged_data <- merged_data %>%
  group_by(countrycode) %>%
  mutate(
    lag_hc = lag(hc),                               # Lagged human capital index
    lag_gov_consumption = lag(csh_g),                # Lagged government consumption
    lag_exports = lag(csh_x),                        # Lagged exports
    lag_inflation = lag(inflation),                  # Lagged inflation
    lag_democracy = lag(v2x_polyarchy),               # Lagged democracy score
    lag_efw = lag(efw)                               # Lagged Economic Freedom Score
  )

# Convert columns to numeric where necessary, replacing non-numeric entries (like "..") with NA
merged_data <- merged_data %>%
  mutate(
    lag_inflation = as.numeric(lag_inflation),
    lag_hc = as.numeric(lag_hc),
    lag_gov_consumption = as.numeric(lag_gov_consumption),
    lag_exports = as.numeric(lag_exports),
    lag_democracy = as.numeric(lag_democracy),
    lag_efw = as.numeric(lag_efw)
  )

# List of columns where missing values need to be imputed
columns_to_impute <- c("treatment_5yr", "lag_hc", "lag_gov_consumption", "lag_exports", "lag_inflation", "lag_democracy", "lag_efw")


# Ensure non-numeric values like ".." are treated as NA
merged_data <- merged_data %>%
  mutate(across(all_of(columns_to_impute), ~ ifelse(.x == "..", NA, as.numeric(.x))))  # Convert non-numeric to NA

# This will ensure that any non-numeric value (such as "..") in the selected columns is replaced with NA.

# Filter out rows where treatment_5yr is NA
merged_data <- merged_data %>%
  filter(!is.na(treatment_5yr))

# Optionally, you may also want to remove rows with missing covariate values
merged_data <- merged_data %>%
  drop_na(lag_hc, lag_gov_consumption, lag_exports, lag_inflation, lag_democracy, lag_efw)

# Step 7: Propensity Score Matching using logistic regression (nearest neighbor)
psm_model <- matchit(treatment_5yr ~ lag_hc + lag_gov_consumption + lag_exports + lag_inflation + lag_democracy + lag_efw,
                     data = merged_data, method = "nearest", distance = "logit")


# Step 8: Check balance of the matching
summary(psm_model)

# Save Propensity Score Matching summary table
save_table_results(summary(psm_model), "Propensity Score Matching Summary",
                   file.path(tables_dir, "PSM_Summary.tex"))

# Extract balance summary for matched data
psm_balance <- summary(psm_model)$sum.matched

# Convert the balance summary to a data frame for saving
psm_balance_df <- as.data.frame(psm_balance)

# Save as a LaTeX table
save_table_results(psm_balance_df, "Propensity Score Matching Summary",
                   file.path(tables_dir, "PSM_Summary.tex"))


# Step 9: Extract matched data
matched_data_psm <- match.data(psm_model)

# Bootstrapping for standard errors (PSM)
bootstrap_psm <- function(data, indices) {
  d <- data[indices, ]
  treated <- d %>% filter(treatment_5yr == 1)
  control <- d %>% filter(treatment_5yr == 0)
  mean(treated$rgdpe_per_capita - control$rgdpe_per_capita)
}
# Perform 250 replications for bootstrapping
psm_se <- boot(data = matched_data_psm, statistic = bootstrap_psm, R = 500)
print(summary(psm_se))  # Bootstrapped standard errors for ATET

# Extract the original estimate (t0) and the bootstrapped replications (t)
boot_estimate <- psm_se$t0         # The original estimate
boot_replicates <- psm_se$t        # The bootstrapped replicates

# Convert them to a data frame for easier saving
boot_results_df <- data.frame(estimate = boot_estimate, replicates = boot_replicates)

# Save the results as a LaTeX table
save_table_results(boot_results_df, "Bootstrapped ATET for PSM",
                   file.path(tables_dir, "PSM_Bootstrap_Results.tex"))


# Step 10: Mahalanobis Nearest Neighbor Matching
mahal_model <- matchit(treatment_5yr ~ lag_hc + lag_gov_consumption + lag_exports + lag_inflation + lag_democracy, + lag_efw,
                       data = merged_data, method = "nearest", distance = "mahalanobis")

# Step 11: Check balance of the Mahalanobis matching
summary(mahal_model)

# Extract balance summary for matched data
mahal_balance <- summary(mahal_model)$sum.matched

# Convert the balance summary to a data frame for saving
mahal_balance_df <- as.data.frame(mahal_balance)

# Save as a LaTeX table
save_table_results(mahal_balance_df, "Mahalanobis Matching Summary",
                   file.path(tables_dir, "MDM_Summary.tex"))

# Step 12: Extract matched data
matched_data_mahal <- match.data(mahal_model)

# Abadie-Imbens (2011) bias-corrected standard errors (MDM)
mdm_se <- Match(Y = merged_data$rgdpe_per_capita, Tr = merged_data$treatment_5yr,
                X = merged_data[, c("lag_hc", "lag_gov_consumption", "lag_exports", "lag_inflation", "lag_democracy", "lag_efw")],
                BiasAdjust = TRUE)
summary(mdm_se)  # Bias-corrected ATET with standard errors



# Extract the single estimate and associated statistics as scalars
mdm_estimate <- mdm_se$est[1, 1]  # Extract the estimate from the matrix
mdm_se_value <- mdm_se$se[1]      # Extract the bias-corrected standard error
mdm_t_stat <- mdm_estimate / mdm_se_value  # Compute the t-stat manually

# Check if the p-value exists and set it to NA if it's missing
mdm_p_value <- if(!is.null(mdm_se$p.val)) mdm_se$p.val[1] else NA

# Create a data frame with these values
mdm_results <- data.frame(
  Estimate = mdm_estimate,
  AI_SE = mdm_se_value,
  T_stat = mdm_t_stat,
  p_value = mdm_p_value
)

# Now you can save the table
save_table_results(mdm_results, "Bias-Corrected ATET for MDM",
                   file.path(tables_dir, "MDM_Bias_Corrected_Results.tex"))




# Step 13: Example analysis: Comparing means between treated and control groups
# Specify the directory to save the LaTeX file
tables_dir <- "C:/Users/ammonsj/Dropbox/Apps/Overleaf/Gender and Economic Freedom/Tables"

# Step 13: Example analysis: Comparing means between treated and control groups
t_test_result <- t.test(matched_data_psm$rgdpe_per_capita[matched_data_psm$treatment_5yr == 1],
                        matched_data_psm$rgdpe_per_capita[matched_data_psm$treatment_5yr == 0])

# Convert t-test result to LaTeX format
t_test_latex <- paste0(
  "\\begin{table}[ht]\n",
  "\\centering\n",
  "\\caption{T-test Results for rgdpe_per_capita}\n",
  "\\begin{tabular}{ll}\n",
  "\\hline\n",
  "Statistic & ", round(t_test_result$statistic, 3), " \\\\ \n",
  "P-value & ", round(t_test_result$p.value, 3), " \\\\ \n",
  "Confidence Interval & [", round(t_test_result$conf.int[1], 3), ", ", round(t_test_result$conf.int[2], 3), "] \\\\ \n",
  "Mean of treated group & ", round(mean(matched_data_psm$rgdpe_per_capita[matched_data_psm$treatment_5yr == 1]), 3), " \\\\ \n",
  "Mean of control group & ", round(mean(matched_data_psm$rgdpe_per_capita[matched_data_psm$treatment_5yr == 0]), 3), " \\\\ \n",
  "\\hline\n",
  "\\end{tabular}\n",
  "\\end{table}"
)

# Export the LaTeX table to a .tex file
writeLines(t_test_latex, con = file.path(tables_dir, "T_Test_Results.tex"))


# Step 14: Example regression analysis on matched data
regression_result <- lm(rgdpe_per_capita ~ treatment_5yr + lag_hc +lag_gov_consumption + lag_exports + lag_inflation + lag_democracy + lag_efw, data = matched_data_psm)
summary(regression_result)

# Save Regression results
save_table_results(regression_result, "Regression Results for Real GDP per Capita",
                   file.path(tables_dir, "Regression_Results.tex"))


# Ungroup the data to ensure that countrycode is not included in the summary
merged_data <- merged_data %>% ungroup()

# Select the relevant numeric variables for propensity score matching (excluding countrycode)
variables_for_summary <- merged_data %>%
  select(lag_hc, lag_gov_consumption, lag_exports, lag_inflation, lag_democracy, lag_efw)

# Calculate the summary statistics manually for each variable
summary_stats <- data.frame(
  N = sapply(variables_for_summary, function(x) sum(!is.na(x))),
  Mean = sapply(variables_for_summary, function(x) mean(x, na.rm = TRUE)),
  StDev = sapply(variables_for_summary, function(x) sd(x, na.rm = TRUE)),
  Min = sapply(variables_for_summary, function(x) min(x, na.rm = TRUE)),
  Max = sapply(variables_for_summary, function(x) max(x, na.rm = TRUE))
)

# View the summary statistics
print(summary_stats)

# Convert the summary statistics to a LaTeX-compatible table using xtable
summary_stats_xtable <- xtable(summary_stats, caption = "Summary Statistics for Variables in Propensity Score Matching")

# Export the LaTeX table to a .tex file
print(summary_stats_xtable, file = file.path(tables_dir, "Summary_Statistics_PSM.tex"),
      type = "latex", include.rownames = FALSE)


# Step 8: Check balance of the matching with bal.tab and save the output
bal_tab_psm <- bal.tab(psm_model)

# Save the balance table as a LaTeX table
bal_tab_psm_xtable <- xtable(as.data.frame(bal_tab_psm$Balance), caption = "Balance Table for PSM")

# Export the LaTeX balance table to Overleaf directory
print(bal_tab_psm_xtable, file = file.path(tables_dir, "Balance_Table_PSM.tex"),
      type = "latex", include.rownames = FALSE, caption.placement = "top")


# Summary output to verify
print(bal_tab_psm)


# Step 9: Create love.plot for PSM balance and save it
love_plot_psm <- love.plot(psm_model, stars = "std", var.order = "unadjusted",
                           abs = TRUE, threshold = 0.1, size = 4)

# Save the love plot to the figures directory as a PDF
ggsave(file.path(figures_dir, "Love_Plot_PSM.pdf"), plot = love_plot_psm, width = 8, height = 6)

# Repeat for Mahalanobis Distance Matching (MDM)

# Step 11: Check balance for MDM and save the output
bal_tab_mahal <- bal.tab(mahal_model)

# Save the balance table for MDM as a LaTeX table
bal_tab_mahal_xtable <- xtable(as.data.frame(bal_tab_mahal$Balance), caption = "Balance Table for MDM")

# Export the LaTeX balance table for MDM to Overleaf directory
print(bal_tab_mahal_xtable, file = file.path(tables_dir, "Balance_Table_MDM.tex"),
      type = "latex", include.rownames = FALSE, caption.placement = "top")

# Step 12: Create love.plot for MDM balance and save it
love_plot_mahal <- love.plot(mahal_model, stars = "std", var.order = "unadjusted",
                             abs = TRUE, threshold = 0.1, size = 4)

# Save the love plot for MDM to the figures directory as a PDF
ggsave(file.path(figures_dir, "Love_Plot_MDM.pdf"), plot = love_plot_mahal, width = 8, height = 6)

# Summary output to verify
print(bal_tab_mahal)


# Define your outcome variables
outcome_vars <- c("rgdpe_per_capita", "maternal_mortality", "sanitation_urban", "life_expectancy")

# Function to compute treatment effect and standard error
compute_treatment_effect <- function(matched_data, outcome_var) {
  treated <- matched_data %>% filter(treatment_5yr == 1)
  control <- matched_data %>% filter(treatment_5yr == 0)
  estimate <- mean(treated[[outcome_var]], na.rm = TRUE) - mean(control[[outcome_var]], na.rm = TRUE)

  # Bootstrap standard error
  bootstrap_func <- function(data, indices) {
    d <- data[indices, ]
    treated_boot <- d %>% filter(treatment_5yr == 1)
    control_boot <- d %>% filter(treatment_5yr == 0)
    mean(treated_boot[[outcome_var]], na.rm = TRUE) - mean(control_boot[[outcome_var]], na.rm = TRUE)
  }
  boot_result <- boot(data = matched_data, statistic = bootstrap_func, R = 500)
  std_error <- sd(boot_result$t)

  list(Estimate = estimate, Std_Error = std_error)
}

# Step 1: Perform PSM
psm_model <- matchit(treatment_5yr ~ lag_hc + lag_gov_consumption + lag_exports + lag_inflation + lag_democracy + lag_efw,
                     data = merged_data, method = "nearest", distance = "logit")

# Step 2: Extract balance tables for PSM
psm_balance <- summary(psm_model)
psm_balance_all <- as.data.frame(psm_balance$sum.all)
psm_balance_matched <- as.data.frame(psm_balance$sum.matched)

# Step 3: Save balance tables for PSM
psm_balance_all_xtable <- xtable(psm_balance_all, caption = "PSM Balance Before Matching")
psm_balance_matched_xtable <- xtable(psm_balance_matched, caption = "PSM Balance After Matching")
print(psm_balance_all_xtable, file = file.path(tables_dir, "PSM_Balance_Before_Matching.tex"),
      type = "latex", include.rownames = TRUE, caption.placement = "top")
print(psm_balance_matched_xtable, file = file.path(tables_dir, "PSM_Balance_After_Matching.tex"),
      type = "latex", include.rownames = TRUE, caption.placement = "top")

# Step 4: Extract matched data for PSM
matched_data_psm <- match.data(psm_model)

# Step 5: Compute treatment effects for PSM
psm_results <- lapply(outcome_vars, function(var) {
  result <- compute_treatment_effect(matched_data_psm, var)
  data.frame(Outcome = var, Estimate = result$Estimate, Std_Error = result$Std_Error)
})
psm_results_df <- do.call(rbind, psm_results)

# Step 6: Save treatment effect estimates for PSM
psm_results_xtable <- xtable(psm_results_df, caption = "Treatment Effect Estimates from PSM")
print(psm_results_xtable, file = file.path(tables_dir, "PSM_Treatment_Effect_Estimates.tex"),
      type = "latex", include.rownames = FALSE, caption.placement = "top")

# Step 7: Perform MDM
mahal_model <- matchit(treatment_5yr ~ lag_hc + lag_gov_consumption + lag_exports + lag_inflation + lag_democracy, lag_efw,
                       data = merged_data, method = "nearest", distance = "mahalanobis")

# Step 8: Extract balance tables for MDM
mahal_balance <- summary(mahal_model)
mahal_balance_all <- as.data.frame(mahal_balance$sum.all)
mahal_balance_matched <- as.data.frame(mahal_balance$sum.matched)

# Step 9: Save balance tables for MDM
mahal_balance_all_xtable <- xtable(mahal_balance_all, caption = "MDM Balance Before Matching")
mahal_balance_matched_xtable <- xtable(mahal_balance_matched, caption = "MDM Balance After Matching")
print(mahal_balance_all_xtable, file = file.path(tables_dir, "MDM_Balance_Before_Matching.tex"),
      type = "latex", include.rownames = TRUE, caption.placement = "top")
print(mahal_balance_matched_xtable, file = file.path(tables_dir, "MDM_Balance_After_Matching.tex"),
      type = "latex", include.rownames = TRUE, caption.placement = "top")

# Step 10: Extract matched data for MDM
matched_data_mahal <- match.data(mahal_model)

# Step 11: Compute treatment effects for MDM
mdm_results <- lapply(outcome_vars, function(var) {
  result <- compute_treatment_effect(matched_data_mahal, var)
  data.frame(Outcome = var, Estimate = result$Estimate, Std_Error = result$Std_Error)
})
mdm_results_df <- do.call(rbind, mdm_results)

# Step 12: Save treatment effect estimates for MDM
mdm_results_xtable <- xtable(mdm_results_df, caption = "Treatment Effect Estimates from MDM")
print(mdm_results_xtable, file = file.path(tables_dir, "MDM_Treatment_Effect_Estimates.tex"),
      type = "latex", include.rownames = FALSE, caption.placement = "top")

# Optional: Save Regression results using stargazer
# Assuming regression_result is an lm or similar object
stargazer(regression_result, type = "latex",
          out = file.path(tables_dir, "Regression_Results 2.tex"),
          title = "Regression Results for Real GDP per Capita")





# Define the function for matching and calculating effects
match_and_effect <- function(merged_data, indices) {
  # Subset the data according to bootstrap sample indices
  data_sample <- merged_data[indices, ]

  # Propensity score matching
  match_model <- matchit(treatment_5yr ~ lag_hc + lag_gov_consumption + lag_exports + lag_inflation + lag_democracy + lag_efw,
                         method = "nearest", data = data_sample)

  # Analyzing matched data
  matched_data <- match.data(match_model)

  # Calculate mean differences for outcome variables
  effects <- sapply(c("rgdpe_per_capita", "maternal_mortality", "sanitation_urban", "life_expectancy"),
                    function(var) mean(matched_data[[var]][matched_data$treatment_5yr == 1], na.rm = TRUE) -
                      mean(matched_data[[var]][matched_data$treatment_5yr == 0], na.rm = TRUE))

  return(effects)
}

# Set seed for reproducibility
set.seed(123)

# Run the bootstrap with 500 iterations
bootstrap_results <- boot(data = merged_data, statistic = match_and_effect, R = 500)

# Calculate summary statistics for bootstrap results
bootstrap_summary <- apply(bootstrap_results$t, 2, function(x) {
  c(Mean = mean(x),
    SE = sd(x),
    Lower_CI = quantile(x, 0.025),
    Upper_CI = quantile(x, 0.975))
})

# Convert summary statistics to a data frame
bootstrap_summary_df <- as.data.frame(t(bootstrap_summary))

# Assign unique row names to the summary statistics
rownames(bootstrap_summary_df) <- c("RGDPE per Capita", "Maternal Mortality", "Urban Sanitation", "Life Expectancy")

# Ensure that the table directory exists (in case it's missing)
if (!dir.exists(tables_dir)) {
  dir.create(tables_dir, recursive = TRUE)
}

# Generate the kable table with summary statistics for HTML output
bootstrap_table <- kable(bootstrap_summary_df,
                         col.names = c("Mean", "Standard Error", "2.5% CI", "97.5% CI"),
                         caption = "Bootstrap Results for Treatment Effects on Various Outcomes") %>%
  kable_styling(bootstrap_options = c("striped", "hover"), full_width = F) %>%
  add_header_above(c(" " = 1, "Outcome Variables" = 4))  # Adding space for row names

# Generate LaTeX formatted table
latex_table <- kable(bootstrap_summary_df, format = "latex",
                     col.names = c("Mean", "Standard Error", "2.5% CI", "97.5% CI"),
                     caption = "Bootstrap Results for Treatment Effects on Various Outcomes") %>%
  kable_styling(bootstrap_options = c("striped", "hover"), full_width = F) %>%
  add_header_above(c(" " = 1, "Outcome Variables" = 4))  # Adding space for row names

# Define the path for saving the LaTeX table
bootstrap_table_path <- file.path(tables_dir, "Bootstrap_Results_Summary_Table.tex")

# Save the LaTeX table to the specified path
writeLines(latex_table, con = bootstrap_table_path)

# Output the HTML table (if needed)
bootstrap_table

# Define the directory for saving figures
figures_dir <- "C:/Users/ammonsj/Dropbox/Apps/Overleaf/Gender and Economic Freedom/Figures"

# List of outcome variables
outcome_vars <- c("rgdpe_per_capita", "life_expectancy")

# Prepare a container for results
results_list <- list()

# Loop through each outcome variable and run gsynth
for (outcome in outcome_vars) {
  # Run gsynth model with treatment_5yr, min.T0, and two-way fixed effects
  model <- gsynth(formula = as.formula(paste(outcome, "~ treatment_5yr")),
                  data = merged_data,
                  index = c("countrycode", "year"),
                  force = "two-way",
                  CV = TRUE,
                  r = c(0, 5),
                  se = TRUE,
                  inference = "parametric",
                  nboots = 50,
                  parallel = FALSE,
                  min.T0 = 7)  # Set minimum pre-treatment periods

  # Store results
  results_list[[outcome]] <- model

  # Plotting and saving plots
  plot_name <- paste(figures_dir, "/gsynth_plot_", outcome, ".png", sep = "")
  png(filename = plot_name)
  plot(model)
  dev.off()
}

# Optional: Display the model summaries
lapply(results_list, summary)

# Load necessary packages
library(stargazer)
library(plm)
library(gsynth)

# Set up paths for saving output tables
tables_dir <- "C:/Users/ammonsj/Dropbox/Apps/Overleaf/Gender and Economic Freedom/Tables"

# Prepare the dataset you're using (replace NAVCO_Complexity_Final_Clean with your dataset)
data <- merged_data  # Example: Use your merged_data variable from earlier

# List of outcome variables
outcome_vars <- c("rgdpe_per_capita", "maternal_mortality", "sanitation_urban", "life_expectancy")

### Function to run models and output LaTeX table for each outcome variable ###
run_models_and_save_table <- function(outcome, title, label, file_name) {
  # Models for the current outcome variable
  model1 <- lm(as.formula(paste(outcome, "~ treatment_5yr")), data = data)
  model2 <- plm(as.formula(paste(outcome, "~ treatment_5yr")), data = data, index = c("year"), model = "within")
  model3 <- plm(as.formula(paste(outcome, "~ treatment_5yr")), data = data, index = c("countrycode"), model = "within")

  # Create the LaTeX table
  stargazer(model1, model2, model3,
            type = "latex",
            title = title,
            label = label,
            dep.var.labels.include = FALSE,  # Turn off automatic dependent variable labels
            column.labels = c("OLS", "Year FE", "Country FE"),
            covariate.labels = c("Treatment (5yr)", "Intercept"),
            keep.stat = c("n", "rsq"),
            star.cutoffs = c(0.1, 0.05, 0.01),
            out = paste0(tables_dir, file_name))
}

# Run models and save tables for each outcome variable
run_models_and_save_table("rgdpe_per_capita", "Impact of Treatment (5yr) on GDP per capita", "tab:gdp_per_capita", "/regression_table_gdp_per_capita.tex")
run_models_and_save_table("maternal_mortality", "Impact of Treatment (5yr) on Maternal Mortality", "tab:maternal_mortality", "/regression_table_maternal_mortality.tex")
run_models_and_save_table("sanitation_urban", "Impact of Treatment (5yr) on Urban Sanitation", "tab:sanitation_urban", "/regression_table_sanitation_urban.tex")
run_models_and_save_table("life_expectancy", "Impact of Treatment (5yr) on Life Expectancy", "tab:life_expectancy", "/regression_table_life_expectancy_secondary.tex")

# Load necessary packages
library(stargazer)
library(plm)

# Set up paths for saving output tables
tables_dir <- "C:/Users/ammonsj/Dropbox/Apps/Overleaf/Gender and Economic Freedom/Tables"

# Prepare the dataset you're using (replace merged_data with your dataset)
data <- merged_data  # Example: Use your merged_data variable from earlier

# List of outcome variables for each table
quality_of_life_outcomes <- c("rgdpe_per_capita", "maternal_mortality", "sanitation_urban", "life_expectancy")

# Create descriptive column labels for each group of outcomes
quality_of_life_labels <- c("GDP per Capita", "Maternal Mortality", "Urban Sanitation", "Life Expectancy")

### Function to run models and output LaTeX table for multiple outcome variables ###
run_models_and_save_table <- function(outcome_vars, column_labels, title, label, file_name) {
  # Store models for each outcome variable
  models <- list()

  for (outcome in outcome_vars) {
    # Run the model indexed by both year and countrycode
    model <- plm(as.formula(paste(outcome, "~ treatment_5yr")),
                 data = data,
                 index = c("countrycode", "year"),
                 model = "within")
    models[[outcome]] <- model
  }

  # Create the LaTeX table using stargazer
  stargazer(models,
            type = "latex",
            title = title,
            label = label,
            dep.var.labels.include = FALSE,  # Turn off automatic dependent variable labels
            column.labels = column_labels,  # Use descriptive column labels
            covariate.labels = c("Treatment (5yr)"),
            keep.stat = c("n", "rsq"),
            star.cutoffs = c(0.1, 0.05, 0.01),
            notes = "Fixed effects by country and year.",
            notes.align = "l",
            out = paste0(tables_dir, file_name))
}

# Run models and save table for "TWFE: Quality of Life" with descriptive labels
run_models_and_save_table(quality_of_life_outcomes,
                          quality_of_life_labels,
                          "TWFE: Quality of Life",
                          "tab:quality_of_life",
                          "/regression_table_quality_of_life.tex")


# Standard Treatment

# ATE calculation function with standard errors and p-values, filtering out NA/NaN values in outcome variables
get_ate_and_p_weightit <- function(weighted_data, outcome) {
  # Filter out rows with NA or NaN in the outcome variable
  weighted_data <- weighted_data %>% filter(!is.na(.[[outcome]]) & !is.nan(.[[outcome]]))

  treated <- weighted_data %>% filter(treatment_5yr == 1)
  control <- weighted_data %>% filter(treatment_5yr == 0)

  # Calculate ATE, standard error, and p-value
  if (nrow(treated) > 0 && nrow(control) > 0) {
    ate_result <- weighted.mean(treated[[outcome]], treated$psm_custom_weights) -
      weighted.mean(control[[outcome]], control$psm_custom_weights)
    se_result <- sqrt(var(treated[[outcome]]) / nrow(treated) + var(control[[outcome]]) / nrow(control))
    p_value <- t.test(treated[[outcome]], control[[outcome]])$p.value
  } else {
    ate_result <- NA
    se_result <- NA
    p_value <- NA
  }

  return(list(ate = ate_result, se = se_result, p = p_value))
}

# Bootstrapping function to handle missing values and resampling
boot_ate <- function(data, outcome, n_boot = 500) {
  ate_list <- replicate(n_boot, {
    sampled_data <- data %>% sample_frac(replace = TRUE)
    res <- get_ate_and_p_weightit(sampled_data, outcome)
    res$ate
  })
  return(mean(ate_list, na.rm = TRUE))  # Mean of the bootstrapped ATEs
}

# Define outcome variables and matching methods
outcomes <- c("rgdpe_per_capita", "maternal_mortality", "sanitation_urban", "life_expectancy")

matching_methods <- list(
  "PS: Nearest Neighbor" = list(method = "nearest", distance = "logit", k = 1),
  "PS: Nearest 2 Neighbors" = list(method = "nearest", distance = "logit", k = 2),
  "PS: Nearest 3 Neighbors" = list(method = "nearest", distance = "logit", k = 3),
  "PS: Nearest 4 Neighbors" = list(method = "nearest", distance = "logit", k = 4),
  "Mahalanobis: Nearest Neighbor" = list(method = "nearest", distance = "mahalanobis", k = 1),
  "Mahalanobis: Nearest 2 Neighbors" = list(method = "nearest", distance = "mahalanobis", k = 2),
  "Mahalanobis: Nearest 3 Neighbors" = list(method = "nearest", distance = "mahalanobis", k = 3),
  "Mahalanobis: Nearest 4 Neighbors" = list(method = "nearest", distance = "mahalanobis", k = 4)
)

# Initialize an empty data frame for results
final_results_table <- data.frame()

# Main loop for matching and ATE calculation with bootstrapping
for (method_name in names(matching_methods)) {
  method_info <- matching_methods[[method_name]]

  # Perform weighting with WeightIt for kernel matching, otherwise use matchit for nearest neighbor or Mahalanobis matching
  if (method_info$method == "npcbps") {
    # Kernel matching using WeightIt (npcbps method)
    weight_model <- weightit(
      treatment_5yr ~ lag_hc + lag_gov_consumption + lag_exports +
        lag_inflation + lag_democracy + lag_efw,
      data = merged_data, method = method_info$method,
      distance = method_info$distance, estimand = "ATE"
    )

    # Attach weights to the entire dataset
    merged_data$psm_custom_weights <- weight_model$weights
    weighted_data <- merged_data
  } else {
    # Nearest neighbor or Mahalanobis matching using MatchIt
    psm_model <- matchit(
      treatment_5yr ~ lag_hc + lag_gov_consumption + lag_exports +
        lag_inflation + lag_democracy + lag_efw,
      data = merged_data, method = method_info$method,
      distance = method_info$distance, k = method_info$k
    )

    # Extract matched data and ensure custom weights are used
    weighted_data <- match.data(psm_model)
    weighted_data$psm_custom_weights <- weighted_data$weights  # Ensure weights are named consistently
  }

  # Create two rows: one for ATE and one for SE
  ate_row <- data.frame(Method = method_name)
  se_row <- data.frame(Method = "")  # Blank method label for standard error row

  # Calculate ATE, SE, and significance for each outcome, including bootstrapping
  for (outcome in outcomes) {
    res <- get_ate_and_p_weightit(weighted_data, outcome)

    # Add bootstrapped ATE
    boot_ate_result <- boot_ate(weighted_data, outcome)

    # Handle NA cases for p-value and ATE
    if (!is.na(res$p)) {
      significance <- ifelse(res$p < 0.01, "***", ifelse(res$p < 0.05, "**", ifelse(res$p < 0.10, "*", "")))
    } else {
      significance <- ""  # No stars if p-value is NA
    }

    # ATE with significance
    if (!is.na(boot_ate_result)) {
      ate_row[[outcome]] <- paste0(round(boot_ate_result, 3), significance)
    } else {
      ate_row[[outcome]] <- NA  # If the result is NA, don't show stars or value
    }

    # Standard error on the next row, no label
    if (!is.na(res$se)) {
      se_row[[outcome]] <- paste0("(", round(res$se, 3), ")")
    } else {
      se_row[[outcome]] <- "(NA)"  # Handle NA standard errors properly
    }
  }

  # Append the two rows (ATE and SE) to the final results table
  final_results_table <- rbind(final_results_table, ate_row, se_row)
}

# Convert the table to LaTeX without extra \begin{table} and \end{table}
latex_table <- xtable(final_results_table)

# Modify the LaTeX output to include \resizebox{\textwidth}{!} but avoid redundant table environments
latex_code <- print(latex_table, type = "latex", include.rownames = FALSE,
                    floating = FALSE, table.placement = NULL)

# Wrap in the correct LaTeX table environment
latex_code <- paste0("\\begin{table}[ht]\n\\centering\n\\caption{Matching Methods Comparison Treatment 5yr. 0.15} \n\\resizebox{\\textwidth}{!}{%\n", latex_code, "}\n\\end{table}")

# Save the LaTeX table to a file
write(latex_code, file = file.path(tables_dir, "Matching_Methods_Comparison.tex"))

#Next treatment

# ATE calculation function with standard errors and p-values, filtering out NA/NaN values in outcome variables
get_ate_and_p_weightit <- function(weighted_data, outcome) {
  # Filter out rows with NA or NaN in the outcome variable
  weighted_data <- weighted_data %>% filter(!is.na(.[[outcome]]) & !is.nan(.[[outcome]]))

  treated <- weighted_data %>% filter(treatment_5yr.1 == 1)
  control <- weighted_data %>% filter(treatment_5yr.1 == 0)

  # Calculate ATE, standard error, and p-value
  if (nrow(treated) > 0 && nrow(control) > 0) {
    ate_result <- weighted.mean(treated[[outcome]], treated$psm_custom_weights) -
      weighted.mean(control[[outcome]], control$psm_custom_weights)
    se_result <- sqrt(var(treated[[outcome]]) / nrow(treated) + var(control[[outcome]]) / nrow(control))
    p_value <- t.test(treated[[outcome]], control[[outcome]])$p.value
  } else {
    ate_result <- NA
    se_result <- NA
    p_value <- NA
  }

  return(list(ate = ate_result, se = se_result, p = p_value))
}

# Bootstrapping function to handle missing values and resampling
boot_ate <- function(data, outcome, n_boot = 500) {
  ate_list <- replicate(n_boot, {
    sampled_data <- data %>% sample_frac(replace = TRUE)
    res <- get_ate_and_p_weightit(sampled_data, outcome)
    res$ate
  })
  return(mean(ate_list, na.rm = TRUE))  # Mean of the bootstrapped ATEs
}

# Define outcome variables and matching methods
outcomes <- c("rgdpe_per_capita", "maternal_mortality", "sanitation_urban", "life_expectancy")

matching_methods <- list(
  "PS: Nearest Neighbor" = list(method = "nearest", distance = "logit", k = 1),
  "PS: Nearest 2 Neighbors" = list(method = "nearest", distance = "logit", k = 2),
  "PS: Nearest 3 Neighbors" = list(method = "nearest", distance = "logit", k = 3),
  "PS: Nearest 4 Neighbors" = list(method = "nearest", distance = "logit", k = 4),
  "Mahalanobis: Nearest Neighbor" = list(method = "nearest", distance = "mahalanobis", k = 1),
  "Mahalanobis: Nearest 2 Neighbors" = list(method = "nearest", distance = "mahalanobis", k = 2),
  "Mahalanobis: Nearest 3 Neighbors" = list(method = "nearest", distance = "mahalanobis", k = 3),
  "Mahalanobis: Nearest 4 Neighbors" = list(method = "nearest", distance = "mahalanobis", k = 4)
)

# Initialize an empty data frame for results
final_results_table <- data.frame()

# Main loop for matching and ATE calculation with bootstrapping
for (method_name in names(matching_methods)) {
  method_info <- matching_methods[[method_name]]

  # Perform weighting with WeightIt for kernel matching, otherwise use matchit for nearest neighbor or Mahalanobis matching
  if (method_info$method == "npcbps") {
    # Kernel matching using WeightIt (npcbps method)
    weight_model <- weightit(
      treatment_5yr.1 ~ lag_hc + lag_gov_consumption + lag_exports +
        lag_inflation + lag_democracy + lag_efw,
      data = merged_data, method = method_info$method,
      distance = method_info$distance, estimand = "ATE"
    )

    # Attach weights to the entire dataset
    merged_data$psm_custom_weights <- weight_model$weights
    weighted_data <- merged_data
  } else {
    # Nearest neighbor or Mahalanobis matching using MatchIt
    psm_model <- matchit(
      treatment_5yr.1 ~ lag_hc + lag_gov_consumption + lag_exports +
        lag_inflation + lag_democracy + lag_efw,
      data = merged_data, method = method_info$method,
      distance = method_info$distance, k = method_info$k
    )

    # Extract matched data and ensure custom weights are used
    weighted_data <- match.data(psm_model)
    weighted_data$psm_custom_weights <- weighted_data$weights  # Ensure weights are named consistently
  }

  # Create two rows: one for ATE and one for SE
  ate_row <- data.frame(Method = method_name)
  se_row <- data.frame(Method = "")  # Blank method label for standard error row

  # Calculate ATE, SE, and significance for each outcome, including bootstrapping
  for (outcome in outcomes) {
    res <- get_ate_and_p_weightit(weighted_data, outcome)

    # Add bootstrapped ATE
    boot_ate_result <- boot_ate(weighted_data, outcome)

    # Handle NA cases for p-value and ATE
    if (!is.na(res$p)) {
      significance <- ifelse(res$p < 0.01, "***", ifelse(res$p < 0.05, "**", ifelse(res$p < 0.10, "*", "")))
    } else {
      significance <- ""  # No stars if p-value is NA
    }

    # ATE with significance
    if (!is.na(boot_ate_result)) {
      ate_row[[outcome]] <- paste0(round(boot_ate_result, 3), significance)
    } else {
      ate_row[[outcome]] <- NA  # If the result is NA, don't show stars or value
    }

    # Standard error on the next row, no label
    if (!is.na(res$se)) {
      se_row[[outcome]] <- paste0("(", round(res$se, 3), ")")
    } else {
      se_row[[outcome]] <- "(NA)"  # Handle NA standard errors properly
    }
  }

  # Append the two rows (ATE and SE) to the final results table
  final_results_table <- rbind(final_results_table, ate_row, se_row)
}

# Convert the table to LaTeX without extra \begin{table} and \end{table}
latex_table <- xtable(final_results_table)

# Modify the LaTeX output to include \resizebox{\textwidth}{!} but avoid redundant table environments
latex_code <- print(latex_table, type = "latex", include.rownames = FALSE,
                    floating = FALSE, table.placement = NULL)

# Wrap in the correct LaTeX table environment
latex_code <- paste0("\\begin{table}[ht]\n\\centering\n\\caption{Matching Methods Comparison Treatment 5yr 0.1} \n\\resizebox{\\textwidth}{!}{%\n", latex_code, "}\n\\end{table}")

# Save the LaTeX table to a file
write(latex_code, file = file.path(tables_dir, "Matching_Methods_Comparison treatment_5yr.1.tex"))

#Next treatment

# ATE calculation function with standard errors and p-values, filtering out NA/NaN values in outcome variables
get_ate_and_p_weightit <- function(weighted_data, outcome) {
  # Filter out rows with NA or NaN in the outcome variable
  weighted_data <- weighted_data %>% filter(!is.na(.[[outcome]]) & !is.nan(.[[outcome]]))

  treated <- weighted_data %>% filter(treatment_5yr.125 == 1)
  control <- weighted_data %>% filter(treatment_5yr.125 == 0)

  # Calculate ATE, standard error, and p-value
  if (nrow(treated) > 0 && nrow(control) > 0) {
    ate_result <- weighted.mean(treated[[outcome]], treated$psm_custom_weights) -
      weighted.mean(control[[outcome]], control$psm_custom_weights)
    se_result <- sqrt(var(treated[[outcome]]) / nrow(treated) + var(control[[outcome]]) / nrow(control))
    p_value <- t.test(treated[[outcome]], control[[outcome]])$p.value
  } else {
    ate_result <- NA
    se_result <- NA
    p_value <- NA
  }

  return(list(ate = ate_result, se = se_result, p = p_value))
}

# Bootstrapping function to handle missing values and resampling
boot_ate <- function(data, outcome, n_boot = 500) {
  ate_list <- replicate(n_boot, {
    sampled_data <- data %>% sample_frac(replace = TRUE)
    res <- get_ate_and_p_weightit(sampled_data, outcome)
    res$ate
  })
  return(mean(ate_list, na.rm = TRUE))  # Mean of the bootstrapped ATEs
}

# Define outcome variables and matching methods
outcomes <- c("rgdpe_per_capita", "maternal_mortality", "sanitation_urban", "life_expectancy")

matching_methods <- list(
  "PS: Nearest Neighbor" = list(method = "nearest", distance = "logit", k = 1),
  "PS: Nearest 2 Neighbors" = list(method = "nearest", distance = "logit", k = 2),
  "PS: Nearest 3 Neighbors" = list(method = "nearest", distance = "logit", k = 3),
  "PS: Nearest 4 Neighbors" = list(method = "nearest", distance = "logit", k = 4),
  "Mahalanobis: Nearest Neighbor" = list(method = "nearest", distance = "mahalanobis", k = 1),
  "Mahalanobis: Nearest 2 Neighbors" = list(method = "nearest", distance = "mahalanobis", k = 2),
  "Mahalanobis: Nearest 3 Neighbors" = list(method = "nearest", distance = "mahalanobis", k = 3),
  "Mahalanobis: Nearest 4 Neighbors" = list(method = "nearest", distance = "mahalanobis", k = 4)
)

# Initialize an empty data frame for results
final_results_table <- data.frame()

# Main loop for matching and ATE calculation with bootstrapping
for (method_name in names(matching_methods)) {
  method_info <- matching_methods[[method_name]]

  # Perform weighting with WeightIt for kernel matching, otherwise use matchit for nearest neighbor or Mahalanobis matching
  if (method_info$method == "npcbps") {
    # Kernel matching using WeightIt (npcbps method)
    weight_model <- weightit(
      treatment_5yr.125 ~ lag_hc + lag_gov_consumption + lag_exports +
        lag_inflation + lag_democracy + lag_efw,
      data = merged_data, method = method_info$method,
      distance = method_info$distance, estimand = "ATE"
    )

    # Attach weights to the entire dataset
    merged_data$psm_custom_weights <- weight_model$weights
    weighted_data <- merged_data
  } else {
    # Nearest neighbor or Mahalanobis matching using MatchIt
    psm_model <- matchit(
      treatment_5yr.125 ~ lag_hc + lag_gov_consumption + lag_exports +
        lag_inflation + lag_democracy + lag_efw,
      data = merged_data, method = method_info$method,
      distance = method_info$distance, k = method_info$k
    )

    # Extract matched data and ensure custom weights are used
    weighted_data <- match.data(psm_model)
    weighted_data$psm_custom_weights <- weighted_data$weights  # Ensure weights are named consistently
  }

  # Create two rows: one for ATE and one for SE
  ate_row <- data.frame(Method = method_name)
  se_row <- data.frame(Method = "")  # Blank method label for standard error row

  # Calculate ATE, SE, and significance for each outcome, including bootstrapping
  for (outcome in outcomes) {
    res <- get_ate_and_p_weightit(weighted_data, outcome)

    # Add bootstrapped ATE
    boot_ate_result <- boot_ate(weighted_data, outcome)

    # Handle NA cases for p-value and ATE
    if (!is.na(res$p)) {
      significance <- ifelse(res$p < 0.01, "***", ifelse(res$p < 0.05, "**", ifelse(res$p < 0.10, "*", "")))
    } else {
      significance <- ""  # No stars if p-value is NA
    }

    # ATE with significance
    if (!is.na(boot_ate_result)) {
      ate_row[[outcome]] <- paste0(round(boot_ate_result, 3), significance)
    } else {
      ate_row[[outcome]] <- NA  # If the result is NA, don't show stars or value
    }

    # Standard error on the next row, no label
    if (!is.na(res$se)) {
      se_row[[outcome]] <- paste0("(", round(res$se, 3), ")")
    } else {
      se_row[[outcome]] <- "(NA)"  # Handle NA standard errors properly
    }
  }

  # Append the two rows (ATE and SE) to the final results table
  final_results_table <- rbind(final_results_table, ate_row, se_row)
}

# Convert the table to LaTeX without extra \begin{table} and \end{table}
latex_table <- xtable(final_results_table)

# Modify the LaTeX output to include \resizebox{\textwidth}{!} but avoid redundant table environments
latex_code <- print(latex_table, type = "latex", include.rownames = FALSE,
                    floating = FALSE, table.placement = NULL)

# Wrap in the correct LaTeX table environment
latex_code <- paste0("\\begin{table}[ht]\n\\centering\n\\caption{Matching Methods Comparison Treatment 5yr 0.125} \n\\resizebox{\\textwidth}{!}{%\n", latex_code, "}\n\\end{table}")

# Save the LaTeX table to a file
write(latex_code, file = file.path(tables_dir, "Matching_Methods_Comparison treatment_5yr.125.tex"))

# Function to create comprehensive jump tables
create_jump_table <- function(data, threshold, treatment_var) {
  
  # Identify all treatment windows (when treatment switches from 0 to 1)
  jump_table <- data %>%
    arrange(countrycode, year) %>%
    group_by(countrycode) %>%
    mutate(
      # Identify start of treatment episodes (0 to 1 transition)
      treatment_start = (!!sym(treatment_var) == 1 & lag(!!sym(treatment_var), default = 0) == 0)
    ) %>%
    filter(treatment_start == TRUE) %>%
    mutate(
      # For each treatment start, calculate the actual 5-year jump
      First_Year = year - 5,  # The jump period starts 5 years before
      Last_Year = year,        # And ends at the current year
      Jump = Gender_Disparity_Index - lag(Gender_Disparity_Index, 5)
    ) %>%
    filter(Jump >= threshold) %>%  # Ensure jump meets threshold
    arrange(countrycode, First_Year) %>%  # ARRANGE BEFORE SELECT
    mutate(Years = paste(First_Year, "-", Last_Year)) %>%
    select(countrycode, Years, Jump)  # Now select final columns
  
  return(jump_table)
}

# Create jump table for 0.15 threshold
jump_table_015 <- create_jump_table(merged_data, 0.15, "treatment_5yr")

# Create LaTeX table
jump_table_015_xtable <- xtable(jump_table_015, 
                                caption = "Countries, Years, and GDI Jump ($\geq$ 0.15)")

# Save the LaTeX table
print(jump_table_015_xtable, 
      file = file.path(tables_dir, "Jump_Table_015_Complete.tex"),
      type = "latex", 
      include.rownames = FALSE, 
      caption.placement = "top",
      tabular.environment = "longtable", 
      floating = FALSE)

# Print summary
cat(paste("\n0.15 Threshold: Found", nrow(jump_table_015), "treatment episodes across", 
          n_distinct(jump_table_015$countrycode), "countries\n"))

# Create jump table for 0.125 threshold  
jump_table_0125 <- create_jump_table(merged_data, 0.125, "treatment_5yr.125")

jump_table_0125_xtable <- xtable(jump_table_0125, 
                                 caption = "Countries, Years, and GDI Jump ($\geq$ 0.125)")

print(jump_table_0125_xtable, 
      file = file.path(tables_dir, "Jump_Table_0125_Complete.tex"),
      type = "latex", 
      include.rownames = FALSE, 
      caption.placement = "top",
      tabular.environment = "longtable", 
      floating = FALSE)

cat(paste("0.125 Threshold: Found", nrow(jump_table_0125), "treatment episodes across", 
          n_distinct(jump_table_0125$countrycode), "countries\n"))

# Create jump table for 0.1 threshold
jump_table_01 <- create_jump_table(merged_data, 0.1, "treatment_5yr.1")

jump_table_01_xtable <- xtable(jump_table_01, 
                               caption = "Countries, Years, and GDI Jump ($\geq$ 0.1)")

print(jump_table_01_xtable, 
      file = file.path(tables_dir, "Jump_Table_01_Complete.tex"),
      type = "latex", 
      include.rownames = FALSE, 
      caption.placement = "top",
      tabular.environment = "longtable", 
      floating = FALSE)

cat(paste("0.1 Threshold: Found", nrow(jump_table_01), "treatment episodes across", 
          n_distinct(jump_table_01$countrycode), "countries\n\n"))

# Optional: Create a summary table showing countries with multiple jumps
multiple_jumps_summary <- jump_table_015 %>%
  group_by(countrycode) %>%
  summarise(
    n_jumps = n(),
    periods = paste(Years, collapse = "; ")
  ) %>%
  filter(n_jumps > 1) %>%
  arrange(desc(n_jumps))

if(nrow(multiple_jumps_summary) > 0) {
  cat("Countries with multiple 0.15+ jumps:\n")
  print(multiple_jumps_summary)
  
  # Save this summary table as well
  multiple_jumps_xtable <- xtable(multiple_jumps_summary,
                                  caption = "Countries with Multiple GDI Jumps ($\geq$ 0.15)")
  print(multiple_jumps_xtable,
        file = file.path(tables_dir, "Multiple_Jumps_Summary.tex"),
        type = "latex",
        include.rownames = FALSE,
        caption.placement = "top")
}

# Load necessary libraries
library(ggplot2)

# Set working directories for saving tables and figures
tables_dir <- "C:/Users/ammonsj/Dropbox/Apps/Overleaf/Gender and Economic Freedom/Tables"  # [INSERT]
figures_dir <- "C:/Users/ammonsj/Dropbox/Apps/Overleaf/Gender and Economic Freedom/Figures"  # [INSERT]

# Ensure directories exist
dir.create(tables_dir, showWarnings = FALSE, recursive = TRUE)
dir.create(figures_dir, showWarnings = FALSE, recursive = TRUE)

# List of countries of interest
countries <- c("BEL", "BOL", "BRA", "CHE", "CHL", "CIV", "DOM", "ECU", "ESP", 
               "GTM", "HTI", "IDN", "LSO", "MUS", "NPL", "PER", "PHL", "PRT", 
               "PRY", "RWA", "SLE", "TGO", "ZAF")

# Filter the dataset to include only the specified countries, and all years within 1970 to 2021
subset_data <- efotw[efotw$countrycode %in% countries & efotw$year >= 1970 & efotw$year <= 2021, ]

# Create the plot
p <- ggplot(subset_data, aes(x = year, y = `Gender Disparity Index`, color = countrycode)) +
  geom_line() +
  labs(title = "Gender Disparity Index Over Time (1970 - 2021)", 
       x = "Year", 
       y = "Gender Disparity Index") +
  theme_minimal() +
  theme(legend.title = element_blank()) # Optionally remove legend title

# Save the plot to the specified directory
ggsave(filename = file.path(figures_dir, "Gender_Disparity_Index_Over_Time.png"), plot = p, width = 8, height = 6)

# Function to save tables (as per your original code)
save_table_results <- function(tab, caption, file_name) {
  stargazer(tab,
            type = "latex",
            title = caption,
            out = file_name)
}

#### NEW CODE####

# ===================================================================================
# FIXED CALLAWAY & SANT'ANNA DID ANALYSIS FOR GDI JUMPS
# ===================================================================================

cat("\n===================================================================================\n")
cat("CALLAWAY & SANT'ANNA DID ANALYSIS: Effect of GDI Jumps on Development Outcomes\n")
cat("===================================================================================\n\n")

# Load required library
library(did)

# --- RELOAD AND PREPARE DATA WITHOUT EXCESSIVE FILTERING ---

# Reload the original merged_data before the heavy filtering
merged_data_cs <- merged_data %>%
  ungroup() %>%
  arrange(countrycode, year)

# Create treatment indicators for CS analysis
# For CS, we need to identify the FIRST treatment time for each unit
merged_data_cs <- merged_data_cs %>%
  group_by(countrycode) %>%
  mutate(
    # Recalculate GDI changes
    GDI_change_5yr_015 = Gender_Disparity_Index - lag(Gender_Disparity_Index, 5),
    GDI_change_5yr_0125 = Gender_Disparity_Index - lag(Gender_Disparity_Index, 5),
    GDI_change_5yr_01 = Gender_Disparity_Index - lag(Gender_Disparity_Index, 5),
    
    # Create treatment indicators
    treat_015 = ifelse(GDI_change_5yr_015 >= 0.15, 1, 0),
    treat_0125 = ifelse(GDI_change_5yr_0125 >= 0.125, 1, 0),
    treat_01 = ifelse(GDI_change_5yr_01 >= 0.1, 1, 0)
  ) %>%
  mutate(
    # Replace NAs with 0 for treatment variables
    treat_015 = replace_na(treat_015, 0),
    treat_0125 = replace_na(treat_0125, 0),
    treat_01 = replace_na(treat_01, 0)
  ) %>%
  ungroup()

# Function to prepare CS data with better handling
prepare_cs_data_fixed <- function(data, treatment_var) {
  
  # First identify first treatment year for each country
  first_treat_data <- data %>%
    filter(!!sym(treatment_var) == 1) %>%
    group_by(countrycode) %>%
    summarise(first_treat_year = min(year, na.rm = TRUE), .groups = 'drop')
  
  # Join back and create CS dataset
  cs_data <- data %>%
    left_join(first_treat_data, by = "countrycode") %>%
    mutate(
      first_treat = replace_na(first_treat_year, 0),
      id = as.numeric(factor(countrycode))
    ) %>%
    rename(time = year) %>%
    select(id, countrycode, time, first_treat, all_of(outcome_vars))
  
  return(cs_data)
}

# Define outcomes
outcome_vars <- c("rgdpe_per_capita", "life_expectancy", "maternal_mortality",
                  "sanitation_urban")

outcome_labels <- c(
  rgdpe_per_capita = "Real GDP per Capita",
  life_expectancy = "Life Expectancy",
  maternal_mortality = "Maternal Mortality",
  sanitation_urban = "Urban Sanitation"

)

# --- ENHANCED CS ANALYSIS FUNCTION ---

run_cs_analysis_fixed <- function(data, treatment_var, threshold_label) {
  
  cat(paste("\n===== ANALYZING TREATMENT THRESHOLD:", threshold_label, "=====\n"))
  
  # Prepare data
  cs_data <- prepare_cs_data_fixed(data, treatment_var)
  
  # Check treated units
  n_treated <- n_distinct(cs_data$countrycode[cs_data$first_treat > 0])
  n_control <- n_distinct(cs_data$countrycode[cs_data$first_treat == 0])
  
  cat(paste("Number of treated countries:", n_treated, "\n"))
  cat(paste("Number of control countries:", n_control, "\n\n"))
  
  # Storage for results
  cs_results <- list()
  
  # Run CS for each outcome
  for (outcome in outcome_vars) {
    cat(paste("Processing", outcome_labels[outcome], "..."))
    
    # Prepare analysis data
    analysis_data <- cs_data %>%
      filter(!is.na(!!sym(outcome))) %>%
      mutate(!!sym(outcome) := as.numeric(.data[[outcome]]))
    
    # Check data availability
    n_treated_outcome <- n_distinct(analysis_data$id[analysis_data$first_treat > 0])
    n_control_outcome <- n_distinct(analysis_data$id[analysis_data$first_treat == 0])
    
    if (n_treated_outcome < 2 || n_control_outcome < 2) {
      cat(paste(" SKIPPED (treated:", n_treated_outcome, "control:", n_control_outcome, ")\n"))
      next
    }
    
    # Run Callaway & Sant'Anna
    tryCatch({
      # Estimate group-time ATTs
      att_gt_results <- att_gt(
        yname = outcome,
        tname = "time",
        idname = "id",
        gname = "first_treat",
        data = analysis_data,
        panel = TRUE,
        control_group = "nevertreated",
        anticipation = 0,
        bstrap = TRUE,
        biters = 100,  # Reduced for speed
        clustervars = "id",
        est_method = "dr",
        xformla = NULL
      )
      
      # Aggregate to event study
      es_effects <- aggte(att_gt_results, type = "dynamic", 
                         min_e = -10, max_e = 10, na.rm = TRUE)
      
      # Aggregate to overall effect
      overall_effects <- aggte(att_gt_results, type = "group", na.rm = TRUE)
      
      # Store results
      cs_results[[outcome]] <- list(
        att_gt = att_gt_results,
        es = es_effects,
        overall = overall_effects,
        att = overall_effects$overall.att,
        se = overall_effects$overall.se,
        pval = 2 * pnorm(-abs(overall_effects$overall.att / overall_effects$overall.se))
      )
      
      # Create event study plot
      es_plot <- ggdid(es_effects) +
        labs(
          title = paste("CS Event Study:", outcome_labels[outcome]),
          subtitle = paste("GDI Jump Threshold ≥", threshold_label),
          x = "Years Relative to Treatment",
          y = "Average Treatment Effect"
        ) +
        theme_minimal(base_size = 12) +
        theme(
          plot.title = element_text(face = "bold"),
          plot.subtitle = element_text(face = "italic")
        ) +
        geom_hline(yintercept = 0, linetype = "dashed", color = "red", alpha = 0.5) +
        coord_cartesian(xlim = c(-10, 10))
      
      # Save as both PNG and PDF
      plot_base_name <- paste0("CS_", outcome, "_", gsub("\\.", "", threshold_label), "_event_study")
      
      # High-resolution PNG
      ggsave(file.path(figures_dir, paste0(plot_base_name, ".png")), 
             es_plot, width = 5, height = 3, dpi = 150)
      
      # PDF
      ggsave(file.path(figures_dir, paste0(plot_base_name, ".pdf")), 
             es_plot, width = 10, height = 6)
      
      cat(" DONE (plots saved as PNG and PDF)\n")
      
    }, error = function(e) {
      cat(paste(" ERROR:", e$message, "\n"))
    })
  }
  
  # Create summary table if we have results
  if (length(cs_results) > 0) {
    summary_df <- data.frame(
      Outcome = character(),
      ATT = numeric(),
      SE = numeric(),
      t_stat = numeric(),
      P_value = numeric(),
      Significance = character(),
      stringsAsFactors = FALSE
    )
    
    for (outcome in names(cs_results)) {
      result <- cs_results[[outcome]]
      t_stat <- result$att / result$se
      sig <- ifelse(result$pval < 0.01, "***",
                   ifelse(result$pval < 0.05, "**",
                         ifelse(result$pval < 0.1, "*", "")))
      
      summary_df <- rbind(summary_df, data.frame(
        Outcome = outcome_labels[outcome],
        ATT = round(result$att, 4),
        SE = round(result$se, 4),
        t_stat = round(t_stat, 3),
        P_value = round(result$pval, 4),
        Significance = sig,
        stringsAsFactors = FALSE
      ))
    }
    
    cat("\n--- Summary of Results ---\n")
    print(summary_df)
    
    # Save LaTeX table
    summary_xtable <- xtable(summary_df,
                            caption = paste("Callaway & Sant'Anna Results - GDI Jump >=", threshold_label),
                            label = paste0("tab:cs_results_", gsub("\\.", "", threshold_label)))
    
    print(summary_xtable, 
          file = file.path(tables_dir, paste0("CS_results_", gsub("\\.", "", threshold_label), ".tex")),
          type = "latex", 
          include.rownames = FALSE, 
          caption.placement = "top")
  }
  
  return(cs_results)
}

# --- RUN ANALYSIS AND GENERATE OUTPUTS ---------------------------------------

cat("\n===================================================================================\n")
cat("EXECUTING CS ANALYSIS FOR ALL THRESHOLDS\n")
cat("===================================================================================\n\n")

# Run the primary Callaway & Sant'Anna analysis for each defined threshold
cs_results_015 <- run_cs_analysis_fixed(merged_data_cs, "treat_015", "0.15")
cs_results_0125 <- run_cs_analysis_fixed(merged_data_cs, "treat_0125", "0.125")
cs_results_01 <- run_cs_analysis_fixed(merged_data_cs, "treat_01", "0.1")

# --- PLOTTING FUNCTIONS WITH CUSTOM STYLING ------------------------------------

#' Apply a custom black and bold style to a ggdid plot
#'
#' @param ggdid_plot A plot object created by the ggdid() function.
#' @return A modified ggplot object with black, bolded points and error bars.
style_cs_plot <- function(ggdid_plot) {
  ggdid_plot +
    # Force both pre- and post-treatment colors to be black
    scale_color_manual(values = c("black", "black")) +
    # Remove the now-redundant color legend
    guides(color = "none") +
    # Make the point estimates larger (bolder)
    geom_point(size = 2.5) +
    # Make the error bars thicker and set a fixed width
    geom_errorbar(size = 0.8, width = 0.25)
}

#' Create and save combined robustness plots for each outcome
#'
#' This function iterates through each outcome variable, generates a styled
#' event study plot for each treatment threshold, and saves them as a
#' combined figure.
create_robustness_plots <- function(results_list) {
  library(gridExtra)
  library(grid)
  
  for (outcome in outcome_vars) {
    plots_list <- list()
    
    # Iterate through the named list of results
    for (threshold_name in names(results_list)) {
      results_set <- results_list[[threshold_name]]
      
      if (outcome %in% names(results_set)) {
        # 1. Create the base plot using ggdid
        base_plot <- ggdid(results_set[[outcome]]$es) +
          labs(
            title = paste("Threshold ≥", sub("cs_results_", "", threshold_name)),
            x = "Years",
            y = "ATT"
          ) +
          theme_minimal(base_size = 20) +
          coord_cartesian(xlim = c(-10, 10)) +
          geom_hline(yintercept = 0, linetype = "dashed", color = "black")
          theme(axis.text.x = element_text(face = "bold", size = 25), color = "black") 
        
        # 2. Apply the custom black and bold style
        styled_plot <- style_cs_plot(base_plot)
        
        # 3. Add the styled plot to our list for combining
        plots_list[[threshold_name]] <- styled_plot
      }
    }
    
    # Proceed if we have more than one plot to combine
    if (length(plots_list) > 1) {
      combined_plot <- arrangeGrob(
        grobs = plots_list,
        ncol = length(plots_list),
        top = textGrob(
          paste("CS Robustness Check:", outcome_labels[outcome]),
          gp = gpar(fontsize = 20, fontface = "bold")
        )
      )
      
      # Save the combined plot in both PNG and PDF formats
      plot_basename <- file.path(figures_dir, paste0("CS_combined_", outcome))
      ggsave(paste0(plot_basename, ".png"), combined_plot, 
             width = 10 * length(plots_list), height = 10, dpi = 150)
      ggsave(paste0(plot_basename, ".pdf"), combined_plot, 
             width = 15 * length(plots_list), height = 15)
             
      cat(paste("- Created combined robustness plot for:", outcome_labels[outcome], "\n"))
    }
  }
}

# Combine results into a named list for easier iteration
all_cs_results <- list(
  "0.15" = cs_results_015,
  "0.125" = cs_results_0125,
  "0.10" = cs_results_01
)

# Generate and save the combined, styled plots
cat("\n--- Generating Combined Robustness Plots ---\n")
create_robustness_plots(all_cs_results)


# --- JUMP IDENTIFICATION TABLES ------------------------------------------------

cat("\n--- Generating GDI Jump Identification Tables ---\n")

#' Identifies all 5-year periods where GDI jump exceeds a threshold.
#'
#' @param data The input dataframe (e.g., merged_data_cs).
#' @param threshold The minimum GDI change to be considered a jump.
#' @return A tibble with country, years of the jump, and jump magnitude.
create_jump_table <- function(data, threshold) {
  data %>%
    ungroup() %>%
    arrange(countrycode, year) %>%
    group_by(countrycode) %>%
    mutate(GDI_change_5yr = Gender_Disparity_Index - lag(Gender_Disparity_Index, 5)) %>%
    filter(!is.na(GDI_change_5yr) & GDI_change_5yr >= threshold) %>%
    transmute(
      Country = countrycode,
      Period = paste(year - 5, year, sep = "–"),
      `GDI Jump` = round(GDI_change_5yr, 3)
    ) %>%
    arrange(Country, Period)
}

# Create and save LaTeX tables for each threshold
thresholds <- c(0.15, 0.125, 0.1)
for (thresh in thresholds) {
  jump_table <- create_jump_table(merged_data_cs, thresh)
  
  # Print summary to console
  cat(paste0("- Threshold ≥ ", thresh, ": Found ", nrow(jump_table), 
             " jump periods across ", n_distinct(jump_table$Country), " countries.\n"))
             
  # Save to LaTeX file
  label_str <- gsub("\\.", "", as.character(thresh))
  xtable(jump_table,
         caption = paste("Country-Periods with GDI Jump ≥", thresh),
         label = paste0("tab:jump_table_", label_str)) %>%
    print(file = file.path(tables_dir, paste0("GDI_Jump_Table_", label_str, ".tex")),
          type = "latex",
          include.rownames = FALSE,
          caption.placement = "top",
          tabular.environment = "longtable",
          floating = FALSE)
}


# --- FINAL SUMMARY ------------------------------------------------------------

cat("\n===================================================================================\n")
cat("ENHANCED CS ANALYSIS COMPLETE\n")
cat("===================================================================================\n")
cat("- Individual event study plots saved as high-resolution PNG and PDF.\n")
cat("- Combined robustness plots (black, bolded) saved where applicable.\n")
cat("- Results and jump identification tables saved in LaTeX format.\n")
cat(paste("\nAll outputs saved to:", figures_dir, "and", tables_dir, "\n"))
cat("===================================================================================\n")

# --- EVENT STUDY SETUP ---

# Ensure data is sorted to handle time-series operations correctly
merged_data <- merged_data %>%
  arrange(countrycode, year)

# Define treatment based on a GDI jump of 0.15 or more over 5 years.
# The 'treatment_5yr' variable from your code is used for this.
# Ensure NA values in the change calculation are handled.
merged_data <- merged_data %>%
  mutate(
    treatment_5yr = ifelse(!is.na(GDI_change_5yr) & GDI_change_5yr >= 0.15, 1, 0)
  )

# Identify the first year a country receives the treatment
treatment_events <- merged_data %>%
  filter(treatment_5yr == 1) %>%
  group_by(countrycode) %>%
  summarise(treatment_year = min(year), .groups = 'drop')

# Join the treatment year back to the main dataset
merged_data <- merged_data %>%
  left_join(treatment_events, by = "countrycode")

# Calculate each observation's time relative to the treatment year
merged_data <- merged_data %>%
  mutate(time_to_treatment = ifelse(!is.na(treatment_year),
                                      year - treatment_year,
                                      NA_integer_))

# Calculate the outcome variable: year-over-year change in property rights for women
merged_data <- merged_data %>%
  group_by(countrycode) %>%
  mutate(prop_rights_change = v2clprptyw - lag(v2clprptyw, 1)) %>%
  ungroup()

# --- VISUALIZATION 1: TIME-SERIES LINE PLOT ---

# Summarize the average change in property rights for each year relative to the jump
time_series_data <- merged_data %>%
  filter(!is.na(time_to_treatment) & time_to_treatment >= -20 & time_to_treatment <= 20) %>%
  filter(!is.na(prop_rights_change)) %>%
  group_by(time_to_treatment) %>%
  summarise(
    mean_change = mean(prop_rights_change, na.rm = TRUE),
    se = sd(prop_rights_change, na.rm = TRUE) / sqrt(n()),
    lower_ci = mean_change - 1.96 * se,
    upper_ci = mean_change + 1.96 * se,
    .groups = 'drop'
  )

# Create the detailed time-series figure
time_series_plot <- ggplot(time_series_data, aes(x = time_to_treatment, y = mean_change)) +
  geom_ribbon(aes(ymin = lower_ci, ymax = upper_ci), alpha = 0.2, fill = "steelblue") +
  geom_line(color = "darkblue", linewidth = 1.2) +
  geom_hline(yintercept = 0, linetype = "solid", color = "gray50") +
  geom_vline(xintercept = 0, linetype = "dashed", color = "red", linewidth = 1) +
  geom_vline(xintercept = -5, linetype = "dotted", color = "orange", linewidth = 1) +
  scale_x_continuous(breaks = seq(-20, 20, 5)) +
  theme_minimal() +
  labs(
    title = "Evolution of Property Rights for Women Around GDI Jumps",
    subtitle = "Mean year-over-year change with 95% confidence intervals",
    x = "Years Relative to GDI Jump (≥0.15 increase over 5 years)",
    y = "Mean Annual Change in Property Rights for Women"
  )

# Save the time-series plot
ggsave(file.path(figures_dir, "property_rights_evolution_line_plot.png"),
       time_series_plot, width = 12, height = 7, dpi = 1200)

# Print the plot to the RStudio viewer
print(time_series_plot)


# --- SUMMARY TABLE (CHRONOLOGICALLY ORDERED) ---

# Create a summary data frame with periods ordered chronologically
period_summary <- merged_data %>%
  mutate(period = factor(
    case_when(
      time_to_treatment >= -20 & time_to_treatment < -5  ~ "Pre-Jump (-20 to -5)",
      time_to_treatment >= -5 & time_to_treatment <= 0  ~ "During Jump (-5 to 0)",
      time_to_treatment > 0 & time_to_treatment <= 20   ~ "Post-Jump (0 to +20)",
      TRUE                                              ~ NA_character_
    ),
    levels = c(
      "Pre-Jump (-20 to -5)",
      "During Jump (-5 to 0)",
      "Post-Jump (0 to +20)"
    )
  )) %>%
  filter(!is.na(period) & !is.na(prop_rights_change)) %>%
  group_by(period) %>%
  summarise(
    `Mean Annual Change` = round(mean(prop_rights_change, na.rm = TRUE), 4),
    `Std. Error` = round(sd(prop_rights_change, na.rm = TRUE) / sqrt(n()), 4),
    `95% CI Lower` = round(`Mean Annual Change` - 1.96 * `Std. Error`, 4),
    `95% CI Upper` = round(`Mean Annual Change` + 1.96 * `Std. Error`, 4),
    `N (Country-Years)` = n(),
    .groups = 'drop'
  ) %>%
  arrange(period)

# Save the summary table as a LaTeX file
stargazer(as.data.frame(period_summary),
          type = "latex",
          summary = FALSE,
          title = "Property Rights for Women: Changes Around GDI Jumps",
          label = "tab:property_rights_summary",
          out = file.path(tables_dir, "property_rights_summary_table.tex"),
          notes = "Note: Changes in property rights for women relative to a GDI jump ≥ 0.15.")

# Print the summary table to the console
cat("\n--- Chronologically Ordered Period Summary Table ---\n")
print(period_summary)


# --- VISUALIZATION 2: CHRONOLOGICAL BAR CHART ---

# Create a bar chart from the ordered summary data
summary_bar_chart <- ggplot(period_summary, aes(x = period, y = `Mean Annual Change`)) +
  geom_col(fill = "steelblue", alpha = 0.9) +
  geom_errorbar(
    aes(ymin = `95% CI Lower`, ymax = `95% CI Upper`),
    width = 0.25,
    color = "black",
    linewidth = 0.5
  ) +
  labs(
    title = "Mean Change in Women's Property Rights Around GDI Jumps",
    subtitle = "Periods are shown in chronological order",
    x = "Period Relative to GDI Jump",
    y = "Mean Annual Change in Property Rights"
  ) +
  theme_minimal() +
  theme(
     plot.title = element_text(face = "bold", size = 14),
     axis.title = element_text(size = 12)
  )

# Save the bar chart
ggsave(file.path(figures_dir, "property_rights_summary_barchart.png"),
       summary_bar_chart, width = 8, height = 6, dpi = 1200)

# Print the final bar chart to the RStudio viewer
print(summary_bar_chart)

# ===================================================================================
# ECONOMIC FREEDOM SUMMARY INDEX BY TREATMENT COHORT
# ===================================================================================

cat("\n===================================================================================\n")
cat("GENERATING ECONOMIC FREEDOM SUMMARY INDEX PLOT BY TREATMENT COHORT\n")
cat("===================================================================================\n\n")

# Create function to identify treated countries for each threshold
identify_treated_countries <- function(data, threshold) {
  data %>%
    ungroup() %>%
    arrange(countrycode, year) %>%
    group_by(countrycode) %>%
    mutate(GDI_change_5yr = Gender_Disparity_Index - lag(Gender_Disparity_Index, 5)) %>%
    filter(!is.na(GDI_change_5yr) & GDI_change_5yr >= threshold) %>%
    pull(countrycode) %>%
    unique()
}

# Identify treated countries for each threshold
treated_015 <- identify_treated_countries(merged_data_cs, 0.15)
treated_0125 <- identify_treated_countries(merged_data_cs, 0.125)
treated_01 <- identify_treated_countries(merged_data_cs, 0.1)

cat(paste("Countries with ≥0.15 jump:", length(treated_015), "\n"))
cat(paste("Countries with ≥0.125 jump:", length(treated_0125), "\n"))
cat(paste("Countries with ≥0.1 jump:", length(treated_01), "\n\n"))

# Calculate average Economic Freedom Summary Index by year for each cohort
efi_by_cohort <- merged_data_cs %>%
  filter(!is.na(Economic_Freedom_Summary_Index)) %>%
  mutate(
    cohort_015 = ifelse(countrycode %in% treated_015, 1, 0),
    cohort_0125 = ifelse(countrycode %in% treated_0125, 1, 0),
    cohort_01 = ifelse(countrycode %in% treated_01, 1, 0)
  ) %>%
  group_by(year) %>%
  summarise(
    `Threshold ≥ 0.15` = mean(Economic_Freedom_Summary_Index[cohort_015 == 1], na.rm = TRUE),
    `Threshold ≥ 0.125` = mean(Economic_Freedom_Summary_Index[cohort_0125 == 1], na.rm = TRUE),
    `Threshold ≥ 0.10` = mean(Economic_Freedom_Summary_Index[cohort_01 == 1], na.rm = TRUE)
  ) %>%
  filter(year >= 1970 & year <= 2021)  # Ensure we're within the data range

# Reshape data for plotting
efi_long <- efi_by_cohort %>%
  pivot_longer(
    cols = starts_with("Threshold"),
    names_to = "Cohort",
    values_to = "EFI_Score"
  ) %>%
  filter(!is.na(EFI_Score))

# Create the plot with different line types
efi_plot <- ggplot(efi_long, aes(x = year, y = EFI_Score, group = Cohort, linetype = Cohort)) +
  geom_line(color = "black", size = 1.2) +
  scale_linetype_manual(
    values = c(
      "Threshold ≥ 0.15" = "solid",
      "Threshold ≥ 0.125" = "dashed",
      "Threshold ≥ 0.10" = "dotted"
    )
  ) +
  scale_x_continuous(
    breaks = seq(1970, 2020, by = 5),
    limits = c(1970, 2021)
  ) +
  scale_y_continuous(
    breaks = seq(0, 10, by = 0.5),
    limits = c(4, 8)  # Adjust based on your data range
  ) +
  labs(
    title = "Average Economic Freedom Summary Index by GDI Jump Treatment Cohort",
    subtitle = "Countries grouped by maximum 5-year GDI increase threshold",
    x = "Year",
    y = "Economic Freedom Summary Index (Average)",
    linetype = "Treatment Cohort"
  ) +
  theme_minimal(base_size = 12) +
  theme(
    # Black and white theme
    panel.background = element_rect(fill = "white", color = NA),
    plot.background = element_rect(fill = "white", color = NA),
    panel.grid.major = element_line(color = "grey80", size = 0.5),
    panel.grid.minor = element_line(color = "grey90", size = 0.25),
    # Bold title
    plot.title = element_text(face = "bold", size = 14, hjust = 0.5),
    plot.subtitle = element_text(size = 11, hjust = 0.5, margin = margin(b = 10)),
    # Legend positioning
    legend.position = "bottom",
    legend.title = element_text(face = "bold", size = 11),
    legend.text = element_text(size = 10),
    legend.key.width = unit(2, "cm"),
    # Axis formatting
    axis.title = element_text(face = "bold", size = 11),
    axis.text = element_text(size = 10),
    axis.text.x = element_text(angle = 45, hjust = 1)
  )

# Save the plot in high resolution (both PNG and PDF)
output_filename <- "Economic_Freedom_by_Treatment_Cohort"

# Save as high-resolution PNG (600 DPI for publication quality)
ggsave(
  filename = file.path(figures_dir, paste0(output_filename, ".png")),
  plot = efi_plot,
  width = 10,
  height = 6,
  dpi = 1200,
  bg = "white"
)

# Save as PDF (vector format, infinitely scalable)
ggsave(
  filename = file.path(figures_dir, paste0(output_filename, ".pdf")),
  plot = efi_plot,
  width = 10,
  height = 6,
  device = "pdf"
)

cat(paste("Economic Freedom plot saved to:", figures_dir, "\n"))

# Print summary statistics for verification
cat("\n--- Summary Statistics by Cohort ---\n")
summary_stats <- efi_long %>%
  group_by(Cohort) %>%
  summarise(
    Mean = round(mean(EFI_Score, na.rm = TRUE), 3),
    SD = round(sd(EFI_Score, na.rm = TRUE), 3),
    Min = round(min(EFI_Score, na.rm = TRUE), 3),
    Max = round(max(EFI_Score, na.rm = TRUE), 3),
    N_Years = n()
  )
print(summary_stats)

cat("\n===================================================================================\n")
cat("ECONOMIC FREEDOM PLOT GENERATION COMPLETE\n")
cat("===================================================================================\n")